\section{Introduction}
\label{sec:introduction}


Since the release of ChatGPT in November 2022, large language models (LLMs) have been adopted widely across software engineering (SE) research~\cite{DBLP:journals/tosem/HouZLYWLLLGW24}, yet the reproducibility and replicability of empirical studies involving LLMs remains uncertain.
% GROUNDING DBLP:journals/tosem/HouZLYWLLLGW24: "Large Language Models (LLMs) have significantly impacted numerous domains, including Software Engineering (SE). Many recent publications have explored LLMs applied to various SE tasks."
Recent findings indicate a low reproducibility of SE studies involving LLMs~\cite{DBLP:journals/corr/abs-2510-25506}. \citeauthor{DBLP:conf/ndss/EvertzRNMNSGPSS26} surveyed 72 peer-reviewed papers from leading security and SE venues published in 2023--2024 and found that every paper contained at least one of nine LLM-specific pitfalls they identified, spanning data collection, pre-training, fine-tuning, prompting, and evaluation; only 15.7\% of the pitfalls present in those papers were explicitly discussed~\cite{DBLP:conf/ndss/EvertzRNMNSGPSS26}.
% GROUNDING DBLP:journals/corr/abs-2510-25506: "For none of the five studies, we were able to fully reproduce the results."
% GROUNDING DBLP:conf/ndss/EvertzRNMNSGPSS26: "all 72 peer-reviewed papers published at leading Security and Software Engineering venues between 2023 and 2024, finding that every paper contains at least one pitfall"
Three characteristics of LLMs make this particularly challenging.
First, their inherent non-determinism causes variability across runs~\cite{DBLP:conf/naacl/SongWLL25, DBLP:journals/corr/abs-2602-07150}, and slight changes can lead to substantially different results~\cite{DBLP:conf/nips/AgarwalSCCB21, DBLP:journals/corr/abs-2412-12509}.
% GROUNDING DBLP:conf/naacl/SongWLL25: "given a particular input, the same LLM may generate significantly different outputs under various decoding configurations, a phenomenon known as non-determinism in generation"
% GROUNDING DBLP:journals/corr/abs-2602-07150: "single-run pass@1 estimates vary by 2.2 to 6.0 percentage points depending on which run is selected, with standard deviations exceeding 1.5 percentage points even at temperature 0"
% GROUNDING DBLP:conf/nips/AgarwalSCCB21: "A small number of training runs coupled with high variability in performance of deep RL algorithms often leads to substantial statistical uncertainty in reported point estimates."
% GROUNDING DBLP:journals/corr/abs-2412-12509: "A 2% variation in the LLM's judgment could lead to significant shifts in the ranking, potentially misrepresenting the true performance of the models"
Second, commercial models evolve beyond version identifiers, so reported performance can change over time~\cite{DBLP:journals/corr/abs-2307-09009}.
% GROUNDING DBLP:journals/corr/abs-2307-09009: "We find that the performance and behavior of both GPT-3.5 and GPT-4 can vary greatly over time."
Third, even for ``open'' models, training data and fine-tuning details often remain undisclosed~\cite{Gibney2024}.
% GROUNDING Gibney2024: "Many of the large language models that power chatbots claim to be open, but restrict access to code and training data."
Moreover, prompt formatting choices alone can shift accuracy by up to 76 percentage points~\cite{DBLP:conf/iclr/Sclar0TS24}, and configured parameters such as temperature affect output variability~\cite{renze2024temperature}.
% GROUNDING DBLP:conf/iclr/Sclar0TS24: "several widely used open-source LLMs are extremely sensitive to subtle changes in prompt formatting in few-shot settings, with performance differences of up to 76 accuracy points"
% GROUNDING renze2024temperature: "A higher temperature makes the output more random, thus favoring more 'creative' predictions."
Hence, not reporting these settings directly affects reproducibility.
Traditional open science practice in SE has focused on releasing source code and datasets as a replication package. The empirical artifact analysis literature in SE focuses more on code repositories and data than on upstream artifacts such as requirements specifications and design documents~\cite{DBLP:journals/jss/LiuHHXZJCM24}. With LLMs, code is often generated from prompts, context files, and other runtime inputs to the model, so reporting must shift left to cover these upstream artifacts in addition to the code and data they produce.
% GROUNDING DBLP:journals/jss/LiuHHXZJCM24: "All the top-starred artifacts include executable code, providing either tools, frameworks, or toolkits."

Although the SE research community has developed guidelines for conducting and reporting specific types of empirical studies such as controlled experiments (e.g.,~\cite{DBLP:books/sp/WohlinRHORW24, DBLP:books/sp/08/SSS2008}), their replications (e.g.,~\cite{DBLP:journals/tse/SantosVOJ21}), or empirical studies in general (e.g., the \emph{ACM SIGSOFT Empirical Standards}~\cite{ralph2021empiricalstandardssoftwareengineering}), none of these address the LLM-specific aspects described above.
% GROUNDING DBLP:books/sp/WohlinRHORW24: "Part II then devotes one chapter to each of the five experiment steps: scoping, planning, execution, analysis, and result presentation."
% GROUNDING DBLP:books/sp/08/SSS2008: "focus on the practical knowledge necessary for conducting, reporting, and using empirical methods in software engineering"
% GROUNDING DBLP:journals/tse/SantosVOJ21: "Provide an analysis procedure with a set of embedded guidelines to aggregate software engineering (SE) replication results."
% GROUNDING ralph2021empiricalstandardssoftwareengineering: "There are three kinds of standards. The General Standard applies to all empirical research."
Previously, a position paper highlighted these issues~\cite{DBLP:conf/wsese/0001BFB25}, but there was no comprehensive community-developed guidance for designing and reporting empirical studies involving LLMs in SE.
% GROUNDING DBLP:conf/wsese/0001BFB25: "While providing a comprehensive set of guidelines is beyond the scope of this position paper, we report a first set of guidelines based on a discussion session with other empiricism experts"

Therefore, we present community-developed guidelines for designing and reporting studies involving LLMs in SE research, co-developed by 22 researchers.
After outlining our \scope, we introduce a taxonomy of \studytypes, then present eight \guidelines.
We complement these with an applicability matrix mapping guidelines to study types and a reporting checklist for authors and reviewers.
For each study type and guideline, we identify relevant examples, both within and outside of SE research.
We maintain the study types and guidelines online as a living resource for the community to use and shape (\href{https://llm-guidelines.org/}{llm-guidelines.org}).


\section{Scope of Guidelines}
\label{sec:scope}

\scopeparagraph{Software Engineering as our Target Discipline}

We target SE research because reporting guidelines from other disciplines do not address its specific needs (see \emph{Related Reporting Guidelines} below). In particular, SE research employs a wide variety of empirical methods~\cite{ralph2021empiricalstandardssoftwareengineering, DBLP:books/sp/WohlinRHORW24}. We organize the LLM-involving subset of these methods into a taxonomy of \studytypes\ifpaper~(Section~\ref{sec:study-types})\fi, and each of our \guidelines specifies its applicability per study type.
% GROUNDING ralph2021empiricalstandardssoftwareengineering: "Each study should be reviewed on its own terms, which is why so many different standards are needed."
% GROUNDING DBLP:books/sp/WohlinRHORW24: "The book provides indispensable information regarding empirical studies, particularly for experiments, but also for systematic literature studies, surveys, and case studies."

\scopeparagraph{Focus on Text-Based Use Cases}

While multi-modal foundation models that use or generate images, audio, or video may also support SE research and practice, we focus on textual use cases of LLMs (e.g., in natural language or programming languages). Many of our guidelines, particularly those concerning model reporting, prompt documentation, and reproducibility, are likely applicable to multi-modal settings as well, but we leave a systematic assessment of this applicability to future work.

\scopeparagraph{Focus on Direct Development or Research Support}

While researchers may use LLMs for many peripheral tasks (e.g., proofreading, spell-checking, translation), our guidelines focus on their direct role in empirical research and engineering practice.
For engineers, we focus on the use of LLMs to automate SE tasks, that is, artificial intelligence (AI) for software engineering (AI4SE) (see~\llmsforengineers).
This includes agentic systems that autonomously plan and execute multi-step tasks using LLMs (see~\design).
For researchers, we focus on the use of LLMs to automate empirical research tasks such as data collection, processing, or analysis (see~\llmsforresearcher).
The 2026 ACM Policy on Authorship likewise requires disclosing AI used in the research itself, not AI used only to assist with writing~\cite{ACM2026}.
% GROUNDING ACM2026: "ACM no longer requires the disclosure of information regarding the use of AI (as distinct from AI used in the conduct of the research itself"

\scopeparagraph{Researchers as our Target Audience}

Our guidelines are intended to help SE researchers design, plan, conduct,
and report empirical studies involving LLMs, and to support scholarly peer
review of such studies. Each guideline includes an \emph{Advice for
Reviewers} subsection with targeted assessment suggestions.
Our guidelines focus on what to report and how. They complement but do not replace methodological guidance for designing specific types of empirical studies.

\scopeparagraph{Reporting Locations}

We use \paper to refer to the manuscript PDF, including any appendices it contains. We use \supplementarymaterial to refer to any artifact, replication package, dataset, or other resource external to the manuscript PDF (e.g., hosted on Zenodo or Figshare). What belongs in the main body versus an appendix is a presentation choice for authors and venues; what belongs inside the manuscript versus outside it is the reporting-location decision the guidelines make. When a recommendation does not specify a location, either \paper or \supplementarymaterial is acceptable.

\scopeparagraph{Navigating These Guidelines}

Our guidelines are structured to support different reading strategies depending on the reader's goal.
\emph{Researchers planning a new study} may start with the taxonomy of study types (see \studytypes) to identify which types apply to their planned work, then consult \ifpaper Table~\ref{tab:guideline-matrix} \else the \href{/guidelines/\#guidelines-by-study-type}{applicability matrix} \fi to determine which guidelines are requirements (\must) and which are recommendations (\should) for those study types. Each guideline section opens with a brief summary, allowing readers to quickly assess relevance before reading the full text.
\emph{Researchers writing up results} may start with the \ifpaper checklist in Appendix~\ref{sec:checklist}, \else \href{/checklist/}{checklist}, \fi which organizes actionable items by typical paper sections (Introduction, Research Design and Methods, Results, etc.). \ifpaper Table~\ref{tab:guidelines} provides a quick reference to all eight guidelines, and Table~\ref{tab:rationale-recommendations} maps each guideline's rationale to its recommendations. \fi
\emph{Reviewers} can use the \emph{Advice for Reviewers} subsection at the end of each guideline for targeted guidance on assessing manuscripts. \ifpaper Table~\ref{tab:guideline-matrix} helps reviewers identify which guidelines apply to the study type under review. \fi

\scopeparagraph{Related Reporting Guidelines}

Reporting guidelines have a long tradition in healthcare research, where CONSORT~\cite{Schulz2010} is the canonical standard for randomized clinical trials. Our reporting checklist (\ifpaper Appendix~\ref{sec:checklist} \else \href{/checklist/}{checklist} \fi) follows its structural template. LLM-specific reporting guidelines have appeared more recently. Outside SE, these include \citeauthor{Gallifant2025}'s TRIPOD-LLM for healthcare~\cite{Gallifant2025}, \citeauthor{DBLP:conf/chi/NavarroSA26}'s HCI guidelines~\cite{DBLP:conf/chi/NavarroSA26}, and \citeauthor{Kapoor2024REFORMS}'s REFORMS for ML-based science~\cite{Kapoor2024REFORMS}. Within SE, they include \citeauthor{DBLP:conf/icse/SallouDP24}'s vision paper on validity threats~\cite{DBLP:conf/icse/SallouDP24}, our earlier workshop position paper~\cite{DBLP:conf/wsese/0001BFB25}, and \citeauthor{DBLP:journals/corr/abs-2601-01954}'s prompt-reporting guideline for automated SE~\cite{DBLP:journals/corr/abs-2601-01954}, derived from a literature review of ICSE, FSE, and ASE papers and a survey of 105 program committee members. Beyond LLM-specific work, the NeurIPS reproducibility checklist~\cite{DBLP:journals/jmlr/PineauVSLBdFL21} prescribes per-submission disclosure for ML papers, covering items such as algorithm description, experimental setup, and statistical reporting.
% GROUNDING Schulz2010: "the 2010 version of the CONSORT Statement, which updates the previous reporting guideline based on new methodological evidence and accumulated experience"
% GROUNDING Gallifant2025: "Large language models (LLMs) are rapidly being adopted in healthcare, necessitating standardized reporting guidelines."
% GROUNDING DBLP:conf/chi/NavarroSA26: "we present a set of considerations for authors, reviewers, and HCI communities around reporting and reviewing papers that involve LLM-integrated systems"
% GROUNDING Kapoor2024REFORMS: "REFORMS was developed based on a consensus of 19 researchers across computer science, data science, mathematics, social sciences, and biomedical sciences"
% GROUNDING DBLP:conf/icse/SallouDP24: "researchers must still be careful as numerous intricate factors can influence the outcomes of experiments involving LLMs"
% GROUNDING DBLP:conf/wsese/0001BFB25: "This paper contributes the first set of holistic guidelines for such studies."
% GROUNDING DBLP:journals/corr/abs-2601-01954: "Based on the findings, we derived a structured guideline that distinguishes essential, desirable, and exceptional reporting elements."
% GROUNDING DBLP:journals/jmlr/PineauVSLBdFL21: "authors were obliged to fill it both at the initial paper submission phase (May 2019), and at the final camera-ready phase (October 2019)"

\citeauthor{DBLP:journals/corr/abs-2601-01954} frame their recommendations as complementing ours. The items their 105 surveyed program committee members endorsed as essential align with what \modelversion, \design, and \limitationsmitigations already require. \citeauthor{DBLP:conf/chi/NavarroSA26}'s HCI guidelines, by contrast, balance \enq{the practical realities of authors' time, cost, and page limitations} with scientific concerns and HCI norms, while our \guidelines aim for comprehensive coverage and use \must and \should to express that balance. On prompt reporting, \citeauthor{DBLP:conf/chi/NavarroSA26} argue for selective disclosure based on each prompt's centrality to author claims, whereas \design recommends fuller disclosure with named exceptions for privacy, anonymity, or confidentiality concerns. On technical evaluation, \citeauthor{DBLP:conf/chi/NavarroSA26} recommend a \enq{modest technical evaluation of the LLM component on a dataset of representative inputs} tailored to HCI research, whereas \benchmarksmetrics asks for established benchmarks, traditional baselines, and inferential statistics where applicable. Peripheral LLM uses such as proofreading, spell-checking, and translation are entirely out of scope (see \emph{Focus on Direct Development or Research Support} above).

Our guidelines name \paper or \supplementarymaterial as the reporting location for each recommendation (see \emph{Reporting Locations} above). In summary, our guidelines apply to SE empirical research broadly, organize recommendations around the taxonomy of seven study types introduced in \studytypes with explicit per-study-type applicability, and pair each recommendation with a target reporting location in \paper or \supplementarymaterial.

\section{Methodology}
\label{sec:methodology}

This project was initiated at the 2024 meeting of the \emph{International Software Engineering Research Network} (ISERN) in Barcelona, Spain, where the first and last author organized a session on guidelines for empirical studies involving LLMs.
These discussions led to a position paper at \emph{WSESE 2025}~\citep{DBLP:conf/wsese/0001BFB25}, which was presented at the \emph{\nth{2} Copenhagen Symposium on Human-Centered Software Engineering AI} (CHASEAI 2024), forming a workstream to collaboratively refine and extend the preliminary guidelines.
Our process follows an established tradition of developing reporting guidelines through expert collaboration, as exemplified by CONSORT for clinical trials~\citep{Begg1996} and, more recently, TRIPOD-LLM for LLM-based prediction studies in medicine~\citep{Gallifant2025}.

We held bi-weekly meetings, in which we decided on the structure of the study type taxonomy and the guidelines.
During these discussions, we added a new study type (\synthesis), refined and extended the guideline sections (in particular \modelversion, \toolarchitecture, \prompts) and covered \benchmarksmetrics as well as \limitationsmitigations.

Then, at least two co-authors were assigned to each study type and guideline to refine and extend the content, followed by a peer review and refinement phase.
Afterwards, the first author reviewed all sections and discussed potential changes with the topic leads.
A requirement for co-authorship was that each author contributed or reviewed content and, most importantly, read the complete guidelines and confirmed that they stand behind all recommendations, not only their own contributions.

To select illustrative examples for each study type and guideline, co-author pairs independently searched for relevant papers using academic databases.
This initial pool was extended based on suggestions from other co-authors.
Identified papers were assessed by one author and subsequently reviewed by a second author; papers that did not add new insights beyond those already covered were not included.
Multiple review rounds were then carried out on the guidelines as a whole, during which selected papers were cross-checked by additional authors and, where necessary, further examples were added or replaced.
The examples are intended to be illustrative, not the result of a systematic review; the involvement of multiple authors at different stages helped mitigate individual selection bias.

After completing the internal review, we invited external experts in empirical SE to review the study types and guidelines.
During the review process of the \emph{Empirical Software Engineering} journal, the structure and content were improved based on reviewers' suggestions.
We self-assessed our guidelines against the quality dimensions of the \emph{SIGSOFT Empirical Standards}~\citep{ralph2021empiricalstandardssoftwareengineering}; this assessment is reported in the conclusion.


\section{Study Types}
\label{sec:study-types}

Empirical guidelines are needed for SE studies involving LLMs to ensure the validity and reproducibility of their results.
However, such guidelines must be tailored to different study types that each pose unique challenges.
%Therefore, understanding the classification of these studies is essential for developing appropriate guidelines.
Therefore, we created a taxonomy of study types to contextualize our recommendations.
%Note that our \guidelines refer to the \studytypes but not the other way around.
We use the term \emph{study type} to refer to categories of LLM involvement in empirical research, rather than to research methodologies (e.g., experiments, case studies). A single empirical study may involve multiple such study types.
Each study type section starts with a \textbf{description}, followed by \textbf{examples} from the SE research community and beyond. The \textbf{advantages} and \textbf{challenges} of using LLMs are discussed in a summary subsection at the end of this section, synthesizing cross-cutting themes across all study types.
See Table~\ref{tab:study-types} for an overview of study types.

%% This sub-table of contents is organizational overkill. 
%% I think it's  a good additional since the guidlines are already a wall of text.

\begin{table}[tb]
\caption{Overview of study types.}
\label{tab:study-types}
\begin{tabular}{@{}lll@{}}
\toprule
\textbf{ID} & \textbf{Section} & \textbf{Topic} \\
\midrule
     & 4.1   & \textbf{\hyperref[sec:llms-as-tools-for-software-engineering-researchers]{LLMs as Tools for Software Engineering Researchers}} \\
S1   & 4.1.1 & \quad \hyperref[sec:llms-as-annotators]{LLMs as Annotators} \\
S2   & 4.1.2 & \quad \hyperref[sec:llms-as-judges]{LLMs as Judges} \\
S3   & 4.1.3 & \quad \hyperref[sec:llms-for-synthesis]{LLMs for Synthesis} \\
S4   & 4.1.4 & \quad \hyperref[sec:llms-as-subjects]{LLMs as Subjects} \\
\midrule
     & 4.2   & \textbf{\hyperref[sec:llms-as-tools-for-software-engineers]{LLMs as Tools for Software Engineers}} \\
S5   & 4.2.1 & \quad \hyperref[sec:studying-llm-usage-in-software-engineering]{Studying LLM Usage in Software Engineering} \\
S6   & 4.2.2 & \quad \hyperref[sec:llms-for-new-software-engineering-tools]{LLMs for New Software Engineering Tools} \\
S7   & 4.2.3 & \quad \hyperref[sec:benchmarking-llms-for-software-engineering-tasks]{Benchmarking LLMs for Software Engineering Tasks} \\
\bottomrule
\end{tabular}
\end{table}

\subsection{LLMs as Tools for Software Engineering Researchers}
\label{sec:llms-as-tools-for-software-engineering-researchers}

LLMs can serve as powerful tools to help researchers conduct empirical studies.
They can automate various tasks such as data collection, pre-processing, and analysis.
For example, LLMs can apply predefined coding guides to qualitative datasets (\annotators), assess the quality of software artifacts (\judges), generate summaries of research papers (\synthesis), or simulate human behavior in empirical studies (\subjects). If LLMs can complete these tasks well enough, they will reduce the time and effort required to conduct a study. However, these applications each pose challenges to validity and reproducibility.


\studytypesubsection{LLMs as Annotators}
\label{sec:llms-as-annotators}

In the annotator role, LLMs perform qualitative coding---the annotation of natural language text such as requirements, interview transcripts, or open-ended survey responses---that researchers would otherwise do by hand.

\studytypeparagraph{Description}
Coding is a time-consuming manual process~\cite{DBLP:journals/ase/BanoHZT24}. LLMs can augment this process, suggest new codes, and label artifacts based on a pre-defined coding guide much faster than humans can~\cite{DBLP:conf/chi/HeHDRH24}. This covers both \emph{closed coding}, where the LLM labels artifacts against a predefined coding guide, and \emph{open coding}, where the LLM proposes new codes from the data. In closed coding, LLM labels can be assessed against those of human coders applying the same guide; in open coding, researchers review the resulting codebook for adequacy (e.g., level of abstraction, redundancies between codes).
% GROUNDING DBLP:journals/ase/BanoHZT24: "qualitative research in SE has grappled with challenges like the time-intensive nature of the work, limitations to scalability due to its manual nature"
% GROUNDING DBLP:conf/chi/HeHDRH24: "The design, testing, and execution of annotation tasks took two days. In contrast, the efficiency of GPT-4 was outstanding."
The extent to which, or under what conditions, LLMs can perform these tasks \textit{effectively} remains an open research question~\cite{DBLP:conf/msr/AhmedDTP25}, and whether they should be used at all for reflexive qualitative analysis is itself contested~\cite{jowsey2025reject}. Indeed, measuring their effectiveness is practically and philosophically challenging. From an interpretivist philosophical perspective, one cannot measure the quality of analysis by comparing one (human or machine) analyst's work to another. From a realist perspective, triangulating across multiple human judges, LLM judges, and other data sources (e.g., whether a pull request is marked as having resolved an issue) improves confidence in the findings but does not prove that any one judge is valid.
% GROUNDING DBLP:conf/msr/AhmedDTP25: "we propose model-model agreement as a predictor of whether a given task is suitable for LLMs at all"
% GROUNDING jowsey2025reject: "We hold the position that analytic approaches such as reflexive thematic analysis are human research practices requiring a subjective, positioned, and reflexive researcher and therefore the use of GenAI in such approaches is not methodologically congruent."

\studytypeparagraph{Examples}
\citeauthor{Huang2023Enhancing} used multiple LLMs for joint annotation of mobile application reviews~\cite{Huang2023Enhancing}.
% GROUNDING Huang2023Enhancing: "we propose a novel approach to constructing mobile application review datasets by leveraging multiple large language models (LLMs) for joint annotation"
They used three models of comparable size with an absolute majority voting rule (i.e., a label is only accepted if it receives more than half of the total votes from the models). This approach slightly outperformed the best individual model tested. 
Meanwhile, \citeauthor{DBLP:conf/msr/AhmedDTP25} examined LLMs as annotators in SE research across five datasets, six LLMs, and ten annotation tasks~\cite{DBLP:conf/msr/AhmedDTP25}.
% GROUNDING DBLP:conf/msr/AhmedDTP25: "applying six state-of-the-art LLMs to ten annotation tasks from five datasets created by prior work"
They found that inter-model agreement strongly correlates with human-model agreement; models performed poorly in tasks where humans also frequently disagreed. They proposed to use model confidence scores to identify specific samples that could be safely delegated to LLMs, potentially reducing human annotation effort without compromising inter-rater agreement.


\studytypesubsection{LLMs as Judges}
\label{sec:llms-as-judges}

\studytypeparagraph{Description}

LLMs can rate properties of software artifacts (e.g. assess code readability, adherence to coding standards, or comment quality) or sort multiple solutions by some attribute. These kinds of judgments are distinct from annotating unstructured text (see Section \annotators).

\studytypeparagraph{Example(s)}

\citeauthor{DBLP:conf/re/LubosFTGMEL24}~\cite{DBLP:conf/re/LubosFTGMEL24} leveraged \href{https://www.llama.com/llama2/}{Llama-2} to evaluate the quality of software requirements statements. 
They prompted the LLM with the text below, where the words in braces reflect the study parameters:

\judgesexample

\citeauthor{DBLP:conf/re/LubosFTGMEL24} evaluated the LLM's output against expert human judges and found moderate agreement for simple requirements and poor agreement for more complex requirements. In contrast, \citeauthor{wang2025multimodalrequirementsdatabasedacceptance}~\cite{wang2025multimodalrequirementsdatabasedacceptance} used an LLM to generate acceptance criteria for user stories, and provided a rubric to an LLM to judge the generated acceptance criteria on interpretable scales (0 to 4). 
 

\studytypeparagraph{Advantages}

Like with annotation, LLMs are often faster and cheaper than human judges, \added{so they support larger datasets for judgment studies}. However, the main constraint that varies by model and budget is the input context size, (the number of tokens one can pass into a model). For example, the upper bound of the context that can be passed to OpenAI's \textsf{o1-mini} model is \href{https://help.openai.com/en/articles/9855712-openai-o1-models-faq-chatgpt-enterprise-and-edu}{32k tokens}. Some people argue that LLMs could be more reliable or less biased than human annotators, but no one has provided any compelling evidence to support this claim.    


\studytypeparagraph{Challenges}

Conceptualizing LLM-as-judge as a measurement instrument (e.g., using an LLM to generate quality ratings for code snippets is akin to measuring the quality of the snippet) gives us a language for evaluating LLM-judges. The LLM-judge should exhibit validity, test-retest reliability, inter-rater reliability, minimal random and systematic error, measurement invariance, etc.~\cite{ralph2024teaching} 

When relying on the judgment of LLMs, researchers must build a \textit{reliable} process for generating judgment labels that considers the non-deterministic nature of LLMs and report the intricacies of that process transparently~\cite{DBLP:journals/corr/abs-2412-12509}.
For example, the order of options has been shown to affect LLM outputs in multiple-choice settings~\cite{DBLP:conf/naacl/PezeshkpourH24}. Additionally, LLMs can behave differently when reviewing their own outputs~\cite{NEURIPS2024_7f1f0218}.

Reliability is a necessary but insufficient condition for validity. 
For example, a reliable LLM might be reliably inaccurate and wrong. Experiments using LLMs to ``peer'' review scholarly articles have not gone well~\cite{DBLP:conf/coling/ZhouC024}. LLM judges are susceptible to adversarial manipulation, where semantics-preserving changes (e.g., misleading code comments) may deceive the judge into accepting flawed artifacts~\cite{DBLP:journals/corr/abs-2510-24367}. For tasks with no single, objectively correct answer, it is not clear whether an LLM judge should reflect the majority opinion or reason  independently. 
The underlying statistical framework of an LLM usually pushes outputs towards the most likely (majority) answer, which might not be the best answer. Questions about bias, accuracy, and trust remain~\cite{DBLP:journals/corr/abs-2406-18403}.

Meanwhile, evaluating and judging large numbers of items, for example, to perform fault localization on the thousands of bugs in large open-source projects, has significant clock time, compute time, and carbon footprint.

Regarding systematic errors, LLMs may suffer from well-documented biases and issues with fairness~\cite{Gallegos2024BiasAF} (e.g., in assessing discussions of pull requests). 

Beyond questions of technical capacity, ethical questions remain, particularly if there is some implicit expectation that a human is judging the output. Involving a human in the judgment loop, for example, to contextualize the scoring, is one approach~\cite{panHumanCenteredDesignRecommendations2024}. 
However, the lack of large-scale ground truth datasets for benchmarking LLM performance in judgment studies hinders progress in this area. There is simply no evidence that LLMs can accurately and reliably judge most relevant properties of SE artifacts, or that offering LLM-based support to human judges improves any specific outcomes.  

\studytypesubsection{LLMs for Synthesis}
\label{sec:llms-for-synthesis}

In the synthesis role, LLMs integrate and interpret information from multiple sources to produce higher-level findings such as themes, patterns, or conceptual frameworks. They can also generate synthetic content (e.g., source code, bug-fix pairs, requirements) for downstream training or evaluation.

\studytypeparagraph{Description}

Unlike annotation (see \annotators), which focuses on categorizing or labeling individual data points, synthesis refers to the process of integrating and interpreting information from multiple sources to generate higher-level insights, identify patterns across datasets, and develop conceptual frameworks or theories. LLMs may be able to support synthesis tasks in SE research by processing and distilling information from qualitative data sources. 
Although synthesis in the preceding notion refers to abstraction and interpretation across multiple data sources, the term is sometimes also used to refer to generating synthetic content (e.g., source code, bug-fix pairs, or requirements) that is then used in downstream tasks to train, fine-tune, or evaluate existing models or tools.
In this case, the synthesis is done primarily using the LLM and its training data; the input is limited to basic instructions and examples.

\studytypeparagraph{Examples} 

Published examples of applying LLMs for synthesis in SE remain scarce. However, some recent work in other domains is instructive~\cite{DBLP:journals/ase/BanoHZT24}.
% GROUNDING DBLP:journals/ase/BanoHZT24: "the integration of Large Language Models (LLMs) like ChatGPT into qualitative research in software engineering, much like in other professional domains"
\citeauthor{DBLP:conf/wsese/BarrosANKKNB25} conducted a systematic mapping study on using LLMs for qualitative research~\cite{DBLP:conf/wsese/BarrosANKKNB25}. The included studies span fields such as healthcare, education, and cultural studies, where LLMs supported qualitative methods, such as grounded theory and thematic analysis, by aiding in pattern identification.
% GROUNDING DBLP:conf/wsese/BarrosANKKNB25: "The studies exploring the use of LLMs as a tool to support qualitative data analysis are concentrated in various fields, including healthcare, education, cultural studies, and analysis of technological applications."
In SE, \citeauthor{DBLP:conf/wsese/LecaVSS25} explored how LLMs have been applied for qualitative data analysis (QDA) and proposed general strategies and guidelines for their application~\cite{DBLP:conf/wsese/LecaVSS25}. \citeauthor{DBLP:journals/corr/abs-2511-14528} complemented this perspective by studying the opportunities and limitations of introducing LLM-based support into QDA, and by formulating recommendations for embedding human–AI collaboration across the thematic analysis phases~\cite{DBLP:journals/corr/abs-2511-14528}. Building on these, subsequent work has proposed hybrid frameworks combining LLM support with human-led QDA. \citeauthor{DBLP:journals/corr/abs-2402-01386} designed an LLM-driven multi-agent system that integrates AI with human decision-making to automate qualitative data analysis methods~\cite{DBLP:journals/corr/abs-2402-01386}. Their system generated initial codes, developed themes, and summarized text. Similarly, \citeauthor{DBLP:journals/corr/abs-2510-18456} compared the performance of humans and LLMs in coding, theme development, definition, and refinement, creating guidelines for a hybrid-LLM framework~\cite{DBLP:journals/corr/abs-2510-18456}.
% GROUNDING DBLP:conf/wsese/LecaVSS25: "This study aimed to investigate how LLMs are currently used in qualitative analysis and how they can be used in software engineering research"
% GROUNDING DBLP:journals/corr/abs-2511-14528: "A reflective workshop with 25 ISERN researchers guided participants through structured discussions of LLM-assisted open coding, theme generation, and theme reviewing, using color-coded canvases to document perceived opportunities, limitations, and recommendations."
% GROUNDING DBLP:journals/corr/abs-2402-01386: "we utilized LLMs to automate and expedite the qualitative data analysis processes. We employed a multi-agent model, where each agent was tasked with executing distinct, individual research related activities"
% GROUNDING DBLP:journals/corr/abs-2510-18456: "We systematically compared human and LLM-generated codes and themes between March and July 2025."
Finally, \citeauthor{DBLP:journals/corr/abs-2506-21138}'s work is an example of using LLMs to create synthetic datasets.
They presented an approach to generate synthetic requirements, showing that they \enq{can match or surpass human-authored requirements for specific classification tasks}~\cite{DBLP:journals/corr/abs-2506-21138}.
% GROUNDING DBLP:journals/corr/abs-2506-21138: "our results establish that synthetic requirements can match or surpass human-authored requirements for specific classification tasks"


% Reverted back to "subject", as participant seems to be primarily used for humans: see  https://dictionary.apa.org/subject and https://dictionary.apa.org/participant
\studytypesubsection{LLMs as Subjects}
\label{sec:llms-as-subjects}

In the subject role, LLMs serve as virtual subjects, generating responses or behaviors that an empirical study would otherwise collect from human participants.

\studytypeparagraph{Description}

In empirical studies, data is collected from participants through methods such as surveys, interviews, or controlled experiments.
LLMs can serve as virtual \emph{subjects} by simulating human behavior and interactions. If LLMs can generate responses that approximate those of human participants, they could be valuable for research involving user interactions, collaborative coding environments, and software usability assessments~\cite{ZHAO2025101167}.
% GROUNDING ZHAO2025101167: "These findings highlight the potential of LLMs not only as linguistic tools, but also as simulated human participants for a wide range of experimental and applied purposes."
To achieve this, prompt engineering techniques are widely employed; for instance, the \textit{Personas Pattern}~\cite{DBLP:conf/naacl/KongZCLQSZWD24} involves tailoring LLM responses to align with predefined profiles or roles that emulate specific user archetypes.
% GROUNDING DBLP:conf/naacl/KongZCLQSZWD24: "By adopting the persona, the LLM can generate more natural, in-character responses tailored to that role."
To serve as virtual subjects, generated responses should be indistinguishable from human-produced texts, consistent with the attitudes and sociodemographic information of the conditioning context (e.g., junior vs.\ senior developers), naturally aligned with the form, tone, and content of the simulated scenario, and reflect patterns in relationships between ideas, demographics, and behavior observed in comparable human data~\cite{DBLP:journals/corr/abs-2209-06899}.
% GROUNDING DBLP:journals/corr/abs-2209-06899: "Generated responses proceed naturally from the conditioning context provided, reliably reflecting the form, tone, and content of the context."

\studytypeparagraph{Examples}

\citeauthor{DBLP:journals/ipm/XuSRGPLSH24}~\cite{DBLP:journals/ipm/XuSRGPLSH24} compiled a list of ways LLMs can support social science research, some of which transfer to empirical SE research. For example, LLMs can emulate human responses and behaviors in simulated interviews and focus groups~\cite{DBLP:journals/ase/GerosaTSS24}. Similarly, \citeauthor{DBLP:conf/icse/BanoGH25}~\cite{DBLP:conf/icse/BanoGH25}, investigated biases in LLM-generated candidate profiles in SE recruitment processes. They found biases favoring male candidates, lighter skin tones, and slim physiques, particularly for senior roles. LLMs may be able to simulate end-user feedback and behavior in usability studies, identify usability issues, and offer suggestions for improvement based on predefined user personas.
% GROUNDING DBLP:journals/ipm/XuSRGPLSH24: "AI for social science, where AI is utilized as a powerful tool to enhance various stages of social science research"
% GROUNDING DBLP:journals/ase/GerosaTSS24: "discussing how LLMs can replicate human responses and behaviors in research settings ... including persona-based prompting for interviews, multi-persona dialogue for focus groups"
% GROUNDING DBLP:conf/icse/BanoGH25: "We generated a total of 300 synthetic profiles, 100 gender-based and 50 gender-neutral profiles from each LLM (GPT-4 and Copilot) specific to SE roles"


\subsection{LLMs as Tools for Software Engineers}
\label{sec:llms-as-tools-for-software-engineers}

LLM-based assistants have become a widely adopted tool for software engineers~\cite{stackoverflow2025survey}, supporting them in various tasks such as code generation and debugging.
% GROUNDING stackoverflow2025survey: "84% of respondents are using or planning to use AI tools in their development process, an increase over last year (76%)."
Researchers have studied how software engineers use LLMs (\llmusage), developed new tools that integrate LLMs (\newtools), and benchmarked LLMs for software engineering tasks (\benchmarkingtasks).


\studytypesubsection{Studying LLM Usage in Software Engineering}
\label{sec:studying-llm-usage-in-software-engineering}

Studies of LLM usage examine how software engineers adopt and integrate LLM-based tools into their workflows.

\studytypeparagraph{Description}

LLM usage studies focus on real-world settings, outside the controlled conditions of benchmarks and experiments.
Researchers can observe software engineers' usage of LLM-based tools in the field, or study if and how they adopt such tools, their usage patterns, as well as perceived benefits and challenges.
Surveys, interviews, observational studies, or analysis of usage logs can provide insights into how LLMs are integrated into development processes, how they influence decision making, and what factors affect their acceptance and effectiveness. 
Such studies can inform improvements for existing LLM-based tools, motivate the design of novel tools, or derive best practices for LLM-assisted software engineering.
They can also uncover risks or deficiencies of existing tools.

\studytypeparagraph{Examples}

Based on a convergent mixed-methods study, \citeauthor{DBLP:journals/tosem/Russo24} found that early adoption of generative AI by software engineers is primarily driven by compatibility with existing workflows~\cite{DBLP:journals/tosem/Russo24}.
% GROUNDING DBLP:journals/tosem/Russo24: "The adoption of AI tools in these early integration stages is primarily driven by their compatibility with existing development workflows."
\citeauthor{DBLP:journals/pacmse/KhojahM0N24} investigated the use of ChatGPT (GPT-3.5) by professional software engineers in a week-long observational study~\cite{DBLP:journals/pacmse/KhojahM0N24}.
% GROUNDING DBLP:journals/pacmse/KhojahM0N24: "We conduct an observational study of 24 professional software engineers who have been using ChatGPT over a period of one week in their jobs"
They found that most developers do not use the code generated by ChatGPT directly but instead use the output as a guide to implement their own solutions.
\citeauthor{DBLP:conf/csee/AzanzaPIG24}'s case study found that LLMs could \enq{enhance personalized, instant onboarding support; however, relying on proprietary external LLMs poses significant data privacy risks}~\cite{DBLP:conf/csee/AzanzaPIG24}. 
% GROUNDING DBLP:conf/csee/AzanzaPIG24: "enhance personalized, instant onboarding support; however, relying on proprietary external LLMs poses significant data privacy risks"
 \citeauthor{DBLP:conf/icsa/JahicS24} found that most participants from 15 software companies they surveyed had already adopted AI (especially ChatGPT) for SE tasks, but cited low-quality outputs, copyright issues, and the risk of proprietary code leaks as barriers to adoption~\cite{DBLP:conf/icsa/JahicS24}. 
 % GROUNDING DBLP:conf/icsa/JahicS24: "Ten survey participants already adopted AI for at least some software engineering tasks, while five of them did not do so."
 % GROUNDING DBLP:conf/icsa/JahicS24: "The results show that low quality of produced output, copyright issues, and the effort for checking solutions before adopting them are and will be the main factors hindering adoption of AI in software engineering."
Retrospective studies that analyze data generated while developers use LLMs can provide additional insights into human-LLM interactions.
For example, researchers can employ data mining methods to build large-scale conversation datasets, such as the DevGPT dataset introduced by \citeauthor{DBLP:conf/msr/XiaoTHM24}~\cite{DBLP:conf/msr/XiaoTHM24}.
% GROUNDING DBLP:conf/msr/XiaoTHM24: "This paper introduces DevGPT, a dataset curated to explore how software developers interact with ChatGPT"
Conversations can then be analyzed using quantitative~\cite{DBLP:conf/msr/RabbiCZI24} and qualitative~\cite{DBLP:conf/msr/MohamedPP24} analysis methods.
% GROUNDING DBLP:conf/msr/RabbiCZI24: "In this study, we quantitatively analyze 1,756 AI-written Python code snippets in the DevGPT dataset and evaluate them for quality and security issues."
% GROUNDING DBLP:conf/msr/MohamedPP24: "an empirical study identifying patterns in the conversation that developers have with LLMs"


\studytypesubsection{LLMs for New Software Engineering Tools}
\label{sec:llms-for-new-software-engineering-tools}

In this role, LLMs serve as components of new tools that assist software engineers (e.g., with code comprehension or test generation) or as autonomous agents that perform multi-step tasks on their behalf.

\studytypeparagraph{Description}

New LLM-based tools support software engineers in their daily tasks, such as code comprehension~\cite{DBLP:conf/chi/YanHWH24} and test case generation~\cite{DBLP:journals/tse/SchaferNET24}.
% GROUNDING DBLP:conf/chi/YanHWH24: "Ivie improved understanding of generated code, and was received by programmers as a highly useful, low distraction complement to the programming assistant."
% GROUNDING DBLP:journals/tse/SchaferNET24: "We implement our approach in TestPilot, an adaptive LLM-based test generation tool for JavaScript that automatically generates unit tests for the methods in a given project's API."
One way of integrating LLM-based tools into software engineers' workflows is using GenAI agents.
Unlike traditional LLM-based tools, these agents are capable of acting autonomously and proactively, are often tailored to meet specific user needs (e.g., via context files or domain-specific tools), and can interact with external environments (e.g., file systems, shells, or web APIs)~\cite{wiesinger2025agents,yang2025code}.
% GROUNDING wiesinger2025agents: "Agents are autonomous and can act independently of human intervention... Agents can also be proactive in their approach to reaching their goals."
% GROUNDING yang2025code: "the emergence of autonomous coding agents capable of multi-step problem decomposition"
From an architectural perspective, GenAI agents can be implemented in various ways~\cite{wiesinger2025agents}, but at their core they run a control loop around the LLM (observe $\rightarrow$ inspect $\rightarrow$ choose $\rightarrow$ act)~\cite{raschka2026components}.
% GROUNDING wiesinger2025agents: "there are many such architectures that can be achieved by the mixing and matching of these components"
% GROUNDING raschka2026components: "An agent is a layer on top, which can be understood as a control loop around the model."
% GROUNDING raschka2026components: "the agent is the system that repeatedly calls the model inside an environment"
Each chosen action (e.g., a tool call) produces a result that the agent feeds back into the next iteration, until the task is complete.
\emph{CoALA}~\cite{DBLP:journals/tmlr/SumersYN024} offers a conceptual framework for organizing such agents.
% GROUNDING DBLP:journals/tmlr/SumersYN024: "modular memory components, a structured action space to interact with internal memory and external environments, and a generalized decision-making process to choose actions"
For coding agents specifically, \citeauthor{raschka2026components} identifies building blocks such as the repository context gathered before each call, the constructed prompt and its tool definitions, structured session memory, and delegation to bounded subagents~\cite{raschka2026components}.
% GROUNDING raschka2026components: "In this article, I lay out six of the main building blocks of a coding agent."
Because these architectures vary, researchers can test and compare them to study how design choices affect downstream performance.

\studytypeparagraph{Examples}

\citeauthor{DBLP:conf/chi/YanHWH24} proposed \emph{IVIE}, a tool integrated into the VS Code graphical interface that generates and explains code using LLMs~\cite{DBLP:conf/chi/YanHWH24}.
% GROUNDING DBLP:conf/chi/YanHWH24: "We present an implementation of Ivie that forks VS Code, applying a modern LLM for timely segmentation and explanation of generated code."
The authors focused more on the presentation, providing a user-friendly interface to interact with the LLM. 
\citeauthor{DBLP:journals/tse/SchaferNET24}~\cite{DBLP:journals/tse/SchaferNET24} presented a large-scale empirical evaluation on the effectiveness of LLMs for automated unit test generation.
% GROUNDING DBLP:journals/tse/SchaferNET24: "This paper presents a large-scale empirical evaluation on the effectiveness of LLMs for automated unit test generation without requiring additional training or manual effort."
They presented \emph{TestPilot}, a tool that implements an approach in which the LLM is provided with prompts that include the signature and implementation of a function under test, along with usage examples extracted from the documentation.
\citeauthor{DBLP:conf/icsm/RichardsW24} introduced a preliminary GenAI agent designed to assist developers in understanding source code by incorporating a reasoning component grounded in the theory of mind~\cite{DBLP:conf/icsm/RichardsW24}.
% GROUNDING DBLP:conf/icsm/RichardsW24: "we designed an LLM-based conversational assistant that provides a personalized interaction based on inferred user mental state (e.g., background knowledge and experience)"
\citeauthor{DBLP:conf/icse-seip/TakerngsaksiriPTTZJLCCW25}~\cite{DBLP:conf/icse-seip/TakerngsaksiriPTTZJLCCW25} present \emph{HULA}, a multi-agent system deployed in Atlassian JIRA that lets engineers refine LLM-generated coding plans and source code, and report acceptance and modification rates from real users.
% GROUNDING DBLP:conf/icse-seip/TakerngsaksiriPTTZJLCCW25: "We design the user interface, implement, and deploy HULA into Atlassian JIRA for internal uses"


\studytypesubsection{Benchmarking LLMs for Software Engineering Tasks}
\label{sec:benchmarking-llms-for-software-engineering-tasks}

Benchmarking studies measure LLM performance on standardized SE tasks against reference outputs and shared metrics.

\studytypeparagraph{Description}

In a benchmark, the reference outputs serve as ground truth, and the metrics measure how well an LLM's outputs match them. Typical tasks include code generation, code summarization, code completion, and code repair~\cite{yang2025code}, but also natural language processing tasks such as anaphora resolution (i.e., the task of identifying the referring expression of a word or phrase occurring earlier in the text). %, interesting for subfields such a Requirements Engineering.
% GROUNDING yang2025code: "Modern LLMs have achieved remarkable capabilities across a wide range of code-related tasks, including code completion, translation, repair, and generation"
Metrics may include general metrics for text generation, such as \emph{ROUGE}, \emph{BLEU}, or \emph{METEOR}~\cite{DBLP:journals/tosem/HouZLYWLLLGW24}, or task-specific metrics, such as \emph{CodeBLEU} for code generation. Benchmarking requires high-quality reference datasets.
% GROUNDING DBLP:journals/tosem/HouZLYWLLLGW24: "Additionally, ROUGE/ROUGE-L [...], METEOR [...], EM (Exact Match) [...] are used in specific studies to evaluate the quality and accuracy of generated code or natural language code descriptions."

\studytypeparagraph{Examples}

In SE, benchmarking may include evaluating an LLM's ability to produce accurate and reliable outputs for a given input (usually a task description, possibly accompanied by data obtained from curated real-world projects or from synthetic SE-specific datasets). \emph{RepairBench}~\cite{DBLP:conf/llm4code/SilvaM25}, for example, contains 574 buggy Java methods and their corresponding fixed versions, which can be used to evaluate the performance of LLMs in code repair tasks. It uses the \emph{Plausible@1} metric (i.e., the probability that the first generated patch passes all test cases) and the \emph{AST Match@1} metric (i.e., the probability that the abstract syntax tree of the first generated patch matches the ground truth patch).
% GROUNDING DBLP:conf/llm4code/SilvaM25: "we define a maximum of 0.2 USD per evaluated bug, or approx. $115.1 for a total of 574 bugs"
\emph{SWE-Bench}~\cite{DBLP:conf/iclr/JimenezYWYPPN24} is a more generic benchmark that contains 2,294 SE Python tasks extracted from GitHub pull requests.
% GROUNDING DBLP:conf/iclr/JimenezYWYPPN24: "We therefore introduce SWE-bench, an evaluation framework including 2294 software engineering problems drawn from real GitHub issues and corresponding pull requests across 12 popular Python repositories."
To score the LLM's task performance, the benchmark validates whether the generated patch successfully compiles and calculates the percentage of passed test cases. Meanwhile, \emph{HumanEval}~\cite{DBLP:journals/corr/abs-2107-03374} is often used to assess code generation. 
% GROUNDING DBLP:journals/corr/abs-2107-03374: "On HumanEval, a new evaluation set we release to measure functional correctness for synthesizing programs from docstrings"


\subsection{Advantages and Challenges}
\label{sec:advantages-challenges}

Using LLMs in empirical SE research, whether as tools for researchers or for software engineers, offers several potential advantages but also raises fundamental challenges that cut across the study types described above.

\studytypeparagraph{Advantages}

The primary advantage of using LLMs for research tasks is \emph{speed, cost reduction, and scalability}. LLMs can annotate, judge, synthesize, and simulate faster and at lower cost than human researchers, with studies showing cost reductions of 50--96\% on various natural language tasks~\cite{DBLP:conf/emnlp/WangLXZZ21} (e.g., \citeauthor{DBLP:conf/chi/HeHDRH24} found that GPT-4 annotation required only two days and 122.08 USD compared to several weeks and 4,508 USD for a comparable MTurk pipeline~\cite{DBLP:conf/chi/HeHDRH24}). Similarly, augmenting or replacing human participants with LLM-generated virtual participants would reduce recruitment effort~\cite{DBLP:conf/vl/Madampe0HO24}. This efficiency unlocks scalability: qualitative research traditionally does not scale well to large samples, but LLMs let researchers process more text than human-only coding allows and support larger judgment datasets.
% GROUNDING DBLP:conf/emnlp/WangLXZZ21: "it costs 50% to 96% less to use labels from GPT-3 than using labels from humans"
% GROUNDING DBLP:conf/chi/HeHDRH24: "The main experiment with GPT-4 totaled only $122.08, compared to the $4,508 spent on MTurk"
% GROUNDING DBLP:conf/vl/Madampe0HO24: "our experience of recruiting software professionals as potential participants in our studies has met with fluctuating success ... This cost us more effort and time than anticipated."

LLMs can also \emph{automate} tasks such as coding qualitative data, assessing artifact quality, and generating summaries, reducing cognitive demands and resources required for qualitative research. Some studies suggest that LLMs may improve \emph{consistency}: ChatGPT's accuracy exceeded crowd workers by approximately 25\%, and LLMs can achieve higher inter-rater agreement than crowd workers and trained annotators~\cite{DBLP:journals/corr/abs-2303-15056}. However, both human and LLM judges exhibit systematic biases, such as preferring outputs with fake citations or rich formatting~\cite{DBLP:conf/emnlp/ChenCLJW24}, and no compelling evidence supports claims that LLMs are less biased than human annotators.
% GROUNDING DBLP:journals/corr/abs-2303-15056: "the zero-shot accuracy of ChatGPT exceeds that of crowd-workers by about 25 percentage points on average"
% GROUNDING DBLP:conf/emnlp/ChenCLJW24: "Results show that human and LLM judges are vulnerable to perturbations to various degrees, and that even the cutting-edge judges possess considerable biases."

LLMs could potentially provide \emph{access to otherwise-inaccessible research contexts}. \textit{If} virtual participants' behavior is sufficiently similar to human participants, LLMs could access underrepresented and hard-to-reach populations, strengthen generalizability and inclusiveness, impute missing data (see \synthesis), and enable research that is ethically problematic with real humans (e.g., questions that would force a human to relive past trauma do not harm an LLM).

Studying real-world LLM-based tools allows researchers to \emph{understand the state of practice}, uncovering usage patterns, adoption rates, and contextual factors, and \emph{generate hypotheses} about how LLMs affect developer productivity, collaboration, and decision-making. From an engineering perspective, developing LLM-based tools \emph{requires less task-specific engineering} than traditional SE approaches such as static analysis or symbolic execution, because a single model can handle diverse inputs without language-specific parsers or hand-crafted rules. Good benchmarks provide \emph{standardized evaluation and model comparison}, reducing research effort and supporting open science by providing common ground for sharing data and results. Benchmarks built for specific SE tasks can help identify LLM weaknesses and support optimization and fine-tuning.

\studytypeparagraph{Challenges}

The stochastic nature of LLM responses, where identical prompts may yield different outputs, \emph{undermines replicability} across all study types, complicating interpretation of experimental results and test-retest reliability. More broadly, \emph{reliability} is a persistent concern: Like any measurement instrument, \judges and \annotators should exhibit validity, test-retest reliability, inter-rater reliability, minimal error, and measurement invariance~\cite{DBLP:books/sp/24/RalphKAA24}. However, LLMs show substantial variability depending on the dataset and task~\cite{DBLP:journals/corr/abs-2306-00176, DBLP:journals/corr/abs-2406-18403}. They are sensitive to prompt variations~\cite{DBLP:journals/corr/abs-2304-11085} and option order~\cite{DBLP:conf/naacl/PezeshkpourH24}, can behave differently when reviewing their own outputs~\cite{DBLP:conf/nips/PanicksseryBF24}, and can be unreliable for high-stakes labeling~\cite{DBLP:conf/chi/Wang0RMM24}. Crucially, \emph{reliability does not imply validity}: a reliable LLM might be reliably inaccurate~\cite{DBLP:conf/coling/ZhouC024}. For tasks with no single correct answer, the statistical framework pushes outputs toward the most likely answer, which may not be the best one~\cite{DBLP:journals/corr/abs-2510-01171, DBLP:journals/corr/abs-2310-06452}. See \modelversion and \limitationsmitigations for reporting guidance on replicability and reliability.
% GROUNDING DBLP:books/sp/24/RalphKAA24: "Multi-item scales or multi-metric instruments and statistical measurement models are required to assess the degree to which instruments measure what they are supposed to measure (construct validity)."
% GROUNDING DBLP:journals/corr/abs-2306-00176: "LLM performance for text annotation is promising but highly contingent on both the dataset and the type of annotation task"
% GROUNDING DBLP:journals/corr/abs-2406-18403: "Models are reliable evaluators on some tasks, but overall display substantial variability depending on the property being evaluated."
% GROUNDING DBLP:journals/corr/abs-2304-11085: "even minor wording alterations in prompts or repeating the identical input can lead to varying outputs"
% GROUNDING DBLP:conf/naacl/PezeshkpourH24: "we demonstrate a considerable performance gap of approximately 13% to 75% in LLMs on different benchmarks, when answer options are reordered"
% GROUNDING DBLP:conf/nips/PanicksseryBF24: "One such bias is self-preference, where an LLM evaluator scores its own outputs higher than others' while human annotators consider them of equal quality."
% GROUNDING DBLP:conf/chi/Wang0RMM24: "they are often erroneous for complex or domain-specific tasks and may introduce bias when compared to human annotations"
% GROUNDING DBLP:conf/coling/ZhouC024: "GPT-3.5 can understand human-written reviews, distinguish emotions, and give consistent scores"
% GROUNDING DBLP:journals/corr/abs-2510-01171: "Post-training alignment often reduces LLM diversity, leading to a phenomenon known as mode collapse."
% GROUNDING DBLP:journals/corr/abs-2310-06452: "RLHF significantly reduces output diversity compared to SFT across a variety of measures."

LLMs and LLM-based tools \emph{evolve rapidly}, and so do practitioners' workflows around them. This combination complicates longitudinal comparisons and can quickly render study findings obsolete.
Findings tied to a specific model version may not transfer to later releases even within the same model family, complicating reproducibility and longitudinal comparisons.

The prevalence of \emph{proprietary tools and opaque training data} limits researchers' ability to assess and mitigate biases, and enables benchmark contamination~\cite{DBLP:journals/corr/abs-2410-16186}. LLMs exhibit \emph{bias} in multiple forms: tendencies to overestimate certain labels~\cite{DBLP:journals/corr/abs-2304-10145}, well-documented fairness issues~\cite{DBLP:journals/coling/GallegosRBTKDYZA24}, and the potential to reinforce prejudices. When simulating human participants, LLMs are \enq{likely to both misportray and flatten the representations of demographic groups}~\cite{DBLP:journals/natmi/WangMD25}. See \openllm and \limitationsmitigations for mitigation strategies.
% GROUNDING DBLP:journals/corr/abs-2410-16186: "LLM creators do not always disclose the exact details of the datasets used, and it is plausible that the model is trained on benchmark datasets intentionally or unintentionally"
% GROUNDING DBLP:journals/corr/abs-2304-10145: "for stance detection, ChatGPT may have a tendency to overestimate certain stances, emphasizing the need for human moderation of labels"
% GROUNDING DBLP:journals/coling/GallegosRBTKDYZA24: "Despite this success, these models can learn, perpetuate, and amplify harmful social biases"
% GROUNDING DBLP:journals/natmi/WangMD25: "We argue analytically for why LLMs are likely to both misportray and flatten the representations of demographic groups"

\emph{Evaluation difficulty} is inherent in studying LLM-based tools and benchmarks. Benchmarks may lack construct validity~\cite{DBLP:books/sp/24/RalphKAA24}, usually do not capture the full complexity of software engineering work~\cite{Chandra2025benchmarks}, and may lead to \emph{overconfidence} and overfitting~\cite{DBLP:journals/corr/abs-2412-03597}. LLMs are also susceptible to adversarial manipulation where semantics-preserving changes may deceive them into accepting flawed artifacts~\cite{DBLP:journals/corr/abs-2510-24367}. See \benchmarksmetrics for detailed guidance on benchmark design, including~\citeauthor{cao2025should}'s guidelines for coding task benchmarks~\cite{cao2025should}.
% GROUNDING DBLP:books/sp/24/RalphKAA24: "We rarely see statistical measurement models in repository mining, benchmarking, or other lab-based quantitative research."
% GROUNDING Chandra2025benchmarks: "Does the suite represent the software engineering work, including its complexity and diversity, that we want AI to help with in the first place?"
% GROUNDING DBLP:journals/corr/abs-2412-03597: "models excel at standardized tests while failing to demonstrate genuine language understanding and adaptability"
% GROUNDING DBLP:journals/corr/abs-2510-24367: "These attacks subtly alter software artifacts to deceive the judge into accepting flawed or malicious code."
% GROUNDING cao2025should: "we introduce a comprehensive code benchmark guideline, HOW2BENCH, comprising 55 rigor-oriented checklists, each annotated with its relative threat level to the benchmark validity"

A fundamental challenge is \emph{philosophical and methodological incongruence} with qualitative research. In a recent open letter, 419 experienced qualitative researchers argued \enq{that analytic approaches such as reflexive thematic analysis are human research practices requiring a subjective, positioned, and reflexive researcher and therefore the use of GenAI in such approaches is not methodologically congruent}~\cite{jowsey2025reject}. LLMs lack the capacity for genuinely reflexive qualitative analysis because they operate on statistical prediction without understanding the meaning of the data being analyzed~\cite{DBLP:journals/corr/abs-2306-13298, DBLP:journals/corr/abs-2510-18456}. Effective use requires structured prompts and careful human oversight~\cite{DBLP:conf/wsese/BarrosANKKNB25}. See \humanvalidation and \limitationsmitigations for guidance.
% GROUNDING jowsey2025reject: "analytic approaches such as reflexive thematic analysis are human research practices requiring a subjective, positioned, and reflexive researcher and therefore the use of GenAI in such approaches is not methodologically congruent"
% GROUNDING DBLP:journals/corr/abs-2306-13298: "Despite LLMs' adeptness at generating linguistically coherent responses, they do not genuinely comprehend the meanings, nuances, and deeper implications of words and phrases."
% GROUNDING DBLP:journals/corr/abs-2510-18456: "LLMs can surface patterns, candidate framings, and alternative readings, but they do not make meaning in the epistemological sense; that remains the role of reflexive human interpretation."
% GROUNDING DBLP:conf/wsese/BarrosANKKNB25: "they are highly dependent on well-structured prompts, which, if poorly designed, can lead to inaccurate or incomplete outputs"

There is \emph{insufficient evidence} for the effectiveness of LLMs in most research roles. No compelling evidence exists that LLMs can accurately simulate human participants~\cite{DBLP:journals/corr/abs-2508-06950} or reliably judge most relevant properties of SE artifacts. Gathering such evidence for each specific usage may be quite difficult~\cite{DBLP:journals/ais/HardingDLL24}. LLMs may be better suited for augmenting rather than replacing human researchers~\cite{DBLP:conf/emnlp/WangLXZZ21, DBLP:conf/chi/HeHDRH24}, but even then, limited evidence exists that augmentation increases \emph{effectiveness} as well as \emph{efficiency}. When LLM outputs are incorrect, they can negatively \emph{affect human judgment}~\cite{DBLP:conf/www/HuangKA23a}. See \humanvalidation for mitigation strategies.
% GROUNDING DBLP:journals/corr/abs-2508-06950: "We conclude that LLMs do not simulate human psychology and recommend that psychological researchers should treat LLMs as useful but fundamentally unreliable tools"
% GROUNDING DBLP:journals/ais/HardingDLL24: "In scenarios like this, the informativeness of the LLM's output is impossible to assess without doing further confirmatory work with human participants."
% GROUNDING DBLP:conf/emnlp/WangLXZZ21: "we propose a novel framework of combining pseudo labels from GPT-3 with human labels, which leads to even better performance"
% GROUNDING DBLP:conf/chi/HeHDRH24: "the research focus might shift from 'using AI to support human labelers' to 'using humans to enhance AI labeling.'"
% GROUNDING DBLP:conf/www/HuangKA23a: "In a case when ChatGPT's decision is wrong, such capability might be a danger to misleading lay people and potentially constructing wrong labels."

Majority voting across multiple outputs improves reliability~\cite{DBLP:journals/corr/abs-2304-11085} but increases \emph{cost and environmental impact}. While open models are available, the most capable ones require substantial hardware; relying on \emph{cloud-based APIs} introduces concerns related to data privacy, security, and replicability. See \limitationsmitigations for sustainability considerations.
% GROUNDING DBLP:journals/corr/abs-2304-11085: "majority decisions on repeated classification outputs of the same input led to improvements in consistency"

Field study findings face \emph{generalizability} challenges because outcomes may be highly sensitive to individual differences, usage patterns, goals, and contexts. Field studies must be \enq{dependable}~\cite{Sullivan2011-ub} beyond traditional validity criteria, which complicates methodology; see \limitationsmitigations for a detailed threat taxonomy.
% GROUNDING Sullivan2011-ub: "Dependability is gained though consistency of data, which is evaluated through transparent research steps and research findings."


\section{Guidelines}
\label{sec:guidelines}

A primary goal of our guidelines is to \emph{enable reproducibility and replicability} of empirical SE studies involving LLMs. As repeating LLM-focused research to verify results lies somewhere between the ACM's definitions of reproducibility (different team, same research artifacts) and repeatability (different team, different research artifacts)~\cite{acm2020artifactreview} due to potential model changes and imperfect research artifacts, we follow \citeauthor{DBLP:journals/corr/abs-2510-25506}~\cite{DBLP:journals/corr/abs-2510-25506} and use the terms interchangeably. While previous guidelines regarding open science and empirical studies still apply, LLM-specific characteristics (e.g. inherent non-determinism~\cite{DBLP:conf/naacl/SongWLL25, YuanNondeterminismLLMInference}, opaque and proprietary models) present additional replicability challenges, which, in turn, demand new guidance.  

Each guideline below begins with a brief \tldr summary, followed by its \textbf{rationale}, \textbf{recommendations}, \textbf{examples}, \textbf{benefits}, \textbf{challenges}, links to the relevant \textbf{study types}.
The \emph{rationale} articulates the underlying principle, i.e., why the guideline matters, while the \emph{recommendations} provide concrete, actionable practices.
Table~\ref{tab:rationale-recommendations} provides a compact overview of this mapping.

The guidelines further contain \textbf{advice for reviewers}.
Broadly, reviewers should use our guidelines to help them interrogate the extent to which a manuscript's authors have done what was practically possible to improve reproducibility. We must neither accept research absent reasonable efforts to improve reproducibility, nor reject research for failing to obtain an impossible goal. We borrow this principle of ``reasonable efforts'' from the SIGSOFT Empirical Standards~\cite{ralph2021empiricalstandardssoftwareengineering}, where it applies to methodological rigor more generally. Of course, papers should acknowledge their limitations, but to determine whether these limitations are reasonable, reviewers should ask ``have the authors done what they could to minimize limitations?'' Our guidelines attempt to capture what ``reasonable efforts'' practically means for LLM-based studies.

To distinguish essential criteria from recommendations, our guidelines use two tiers.
A~\must criterion is a \emph{requirement}. Studies that intend to follow these
guidelines are expected to meet all \must criteria.
A~\should criterion represents a desired practice that strengthens a study's rigor or transparency.
However, there may be valid reasons to deviate in particular circumstances (e.g., resource constraints, inapplicability to a specific study context or type).
Nonetheless, studies that deviate from a \should criterion should briefly justify the
deviation and discuss its potential impact (e.g., on validity or reproducibility).
Table~\ref{tab:guidelines} lists our eight guidelines.
Table~\ref{tab:guideline-matrix} provides an overview of how each guideline applies to each study type.
Cells marked \iconM indicate that the guideline is a \must requirement for the study type and cells marked \iconS indicate that the guideline \should be followed.
Cells marked -- indicate that the guideline is typically not applicable or not feasible for the study type.

The following sections indicate which information we expect researchers to report, and whether it should be in the \paper or \supplementarymaterial. Where a publication venue's page limits hinder reporting all expected elements in the \paper, it is better to report essential information in the \supplementarymaterial than not at all.
The \supplementarymaterial should be published according to the \emph{ACM SIGSOFT Open Science Policies}~\citep{daniel_graziotin_2024_10796477}.

\begin{table}[tb]
\caption{Overview of guidelines.}
\label{tab:guidelines}
\begin{tabular}{@{}lll@{}}
\toprule
\textbf{ID} & \textbf{Section} & \textbf{Guideline} \\
\midrule
G1 & \ref*{sec:declare-llm-usage-and-role} & \hyperref[sec:declare-llm-usage-and-role]{Declare LLM Usage and Role} \\
G2 & \ref*{sec:report-model-version-configuration-and-customizations} & \hyperref[sec:report-model-version-configuration-and-customizations]{Report Model Version, Configuration, and Customizations} \\
G3 & \ref*{sec:report-tool-architecture-beyond-models} & \hyperref[sec:report-tool-architecture-beyond-models]{Report Tool Architecture Beyond Models} \\
G4 & \ref*{sec:report-prompts-their-development-and-interaction-logs} & \hyperref[sec:report-prompts-their-development-and-interaction-logs]{Report Prompts, their Development, and Interaction Logs} \\
G5 & \ref*{sec:use-human-validation-for-llm-outputs} & \hyperref[sec:use-human-validation-for-llm-outputs]{Use Human Validation for LLM Outputs} \\
G6 & \ref*{sec:use-an-open-llm-as-a-baseline} & \hyperref[sec:use-an-open-llm-as-a-baseline]{Use an Open LLM as a Baseline} \\
G7 & \ref*{sec:use-suitable-baselines-benchmarks-and-metrics} & \hyperref[sec:use-suitable-baselines-benchmarks-and-metrics]{Use Suitable Baselines, Benchmarks, and Metrics} \\
G8 & \ref*{sec:report-limitations-and-mitigations} & \hyperref[sec:report-limitations-and-mitigations]{Report Limitations and Mitigations} \\
\bottomrule
\end{tabular}
\end{table}

\input{_summary/matrix.tex}

\subsection{Declare LLM Usage and Role (G1)}
\label{sec:declare-llm-usage-and-role}

\vspace{0.5\baselineskip}
\begin{framed}
\tldr: Researchers \must disclose any use of LLMs to support empirical studies in their \paper. They \should report why they used an LLM, what tasks were automated tasks, and the expected benefits of this automation.
\end{framed}

\guidelinesubsubsection{Rationale}

Transparency about LLM involvement is a prerequisite for informed assessment of a study's scope, limitations, and potential biases.
Without explicit disclosure, readers cannot evaluate how the LLM's characteristics may have influenced the research process or its outcomes.

\guidelinesubsubsection{Recommendations}

When conducting any kind of empirical study involving LLMs, researchers \must clearly declare that an LLM was used (see \scope for what we consider relevant research support).
This \should be done in a suitable section of the \paper, for example, in the introduction or research methods section; \citeauthor{cheng2025writing}~\cite{cheng2025writing} argue specifically for the methods section, since acknowledgments are easily missed at the end of a paper.
% GROUNDING cheng2025writing: "the most helpful and transparent means for describing the use of these tools is within the methods section of a manuscript"
The ACM Policy on Authorship requires authors to describe in the methods section any use of AI in the research itself~\cite{ACM2026}.
% GROUNDING ACM2026: "the specific use(s) of AI tools must be described in detail in the methods section of the Work"

Beyond generic declarations, researchers \should report the exact purpose of using an LLM in a study, the tasks it was used to automate, and the expected benefits in the \paper.
A sufficient declaration specifies not only \emph{that} an LLM was used, but also \emph{which} LLM (name and version), \emph{how} it was used (e.g., as an annotator, code generator, or judge), and \emph{where} in the research process it was employed (e.g., data collection, analysis, or synthesis).

When the LLM is central to the study (e.g., as the main tool being evaluated or as a core component of the research method), the declaration \should be prominent and detailed, appearing in the methodology section with cross-references to the specific guidelines that apply (e.g., Sections \modelversion, \design, and \traces).
When the LLM's role is more tangential (e.g., used for a single preprocessing step), a brief but explicit statement in the methodology section is sufficient.
When a study assigns multiple distinct roles to LLMs (e.g., one model generates evaluation data while another scores outputs), each role \should be declared separately.
In each case, the disclosure \must be specific enough for readers to assess how the LLM's involvement may affect the study's validity and reproducibility.

\guidelinesubsubsection{Examples}

The \emph{ACM Policy on Authorship} requires reporting AI-generated artifacts such as code, datasets, and figures where they underlie a study's conclusions~\cite{ACM2026}.
% GROUNDING ACM2026: "the creation of artifacts that are directly relevant to the conclusions of the research, such as code, datasets, and charts or figures that rely on the AI tools"
Reporting in the methods section keeps this disclosure visible under double-blind review, where end-of-paper acknowledgments are often removed.
An example of an LLM disclosure beyond writing support can be found in a recent paper by \citeauthor{DBLP:conf/re/LubosFTGMEL24}~\cite{DBLP:conf/re/LubosFTGMEL24}, in which they write in the methodology section:
% GROUNDING DBLP:conf/re/LubosFTGMEL24: "We conducted an LLM-based evaluation of requirements utilizing the Llama 2 language model with 70 billion parameters, fine-tuned to complete chat responses."

\begin{quote}
\it
``We conducted an LLM-based evaluation of requirements utilizing the Llama 2 language model with 70 billion parameters, fine-tuned to complete chat responses...''
\end{quote}

\noindent A more contemporary declaration could similarly state:

\begin{quote}
\it
``We used Claude Opus 4.7 via the Anthropic API to synthesize themes from interview transcripts, with all prompts and conversation logs published as supplementary material.''
\end{quote}

\citeauthor{DBLP:journals/corr/abs-2601-11895}'s DevBench paper illustrates separate disclosure of multiple LLM roles: \emph{GPT-4o} generates the synthetic benchmark instances, nine models (including \emph{Claude 4 Sonnet}, \emph{GPT-4.1}, \emph{DeepSeek-V3}, and \emph{Ministral-3B}) are the evaluation subjects, and \emph{o3-mini} is the LLM judge that scores completions for relevance and helpfulness~\cite{DBLP:journals/corr/abs-2601-11895}.
% GROUNDING DBLP:journals/corr/abs-2601-11895: "Synthetic instances were generated with OpenAI's GPT-4o, chosen for its fluency, reasoning, and code generation capabilities"

\guidelinesubsubsection{Benefits}

Transparency in the use of LLMs helps other researchers understand the context and scope of the study, supporting interpretation and comparison of the results.
Realizing these benefits requires reporting the LLM's exact role and version (see \modelversion).

\guidelinesubsubsection{Challenges}

Declaring LLM usage requires only a brief statement and no additional experiments, making compliance straightforward.
One challenge might be authors' reluctance to disclose LLM usage for valid use cases, because they fear that AI-generated content makes reviewers think that the authors' work is less original.
In fact, there is evidence suggesting that AI disclosure can negatively affect trust in authors~\cite{SCHILKE2025104405}.
% GROUNDING SCHILKE2025104405: "Thirteen experiments consistently demonstrate that actors who disclose their AI usage are trusted less than those who do not."
However, the \emph{ACM Policy on Authorship} requires disclosure of AI used in the research itself, not AI used only to assist with writing~\cite{ACM2026}.
% GROUNDING ACM2026: "ACM no longer requires the disclosure of information regarding the use of AI (as distinct from AI used in the conduct of the research itself"
Our guidelines focus on such use (see \scope), not on proofreading or writing support.

\guidelinesubsubsection{Study Types}

Researchers \must follow this guideline for all study types.
The specific focus of the declaration varies by study type.
For \annotators, \judges, \synthesis, and \subjects, researchers \must declare the specific role assigned to the LLM (e.g., annotator, judge, synthesizer, or simulated participant).
For \llmusage, researchers \must clarify which LLM(s) the observed participants used and under which conditions.
For \newtools, researchers \must declare the LLM's role within the tool architecture and its contribution to the tool's functionality.
For \benchmarkingtasks, researchers \must declare which LLMs were benchmarked and for which tasks.

\guidelinesubsubsection{Advice for Reviewers}

The most common problem with disclosure is incompleteness or vagueness about how the LLM was used. If the paper says ``we used LLM X to help with task Y'' without specifying how, reviewers should request clarification. Such requests are typically minor revisions unless the missing details may reveal methodological problems.

If undisclosed LLM use is suspected, the reviewer should consult their editor or program chair. When the evidence is conclusive, the key question is the degree to which undisclosed use affects the study's contribution, ranging from negligible (e.g., word choice in a single sentence) to severe (e.g., generating the reported data or results).

\guidelinesubsubsection{See Also}
\begin{itemize}[label=$\rightarrow$]
    \item \seemodelversion: The disclosure is incomplete without naming the specific model and version.
    \item \seedesign: When the LLM lives inside a tool or agent, authors must also describe that tool's architecture and prompts.
    \item \seetraces: Session traces show what the LLM did during the study.
\end{itemize}

\subsection{Report Model Version, Configuration, and Customizations (G2)}
\label{sec:report-model-version-configuration-and-customizations}

\vspace{0.5\baselineskip}
\begin{framed}
\tldr: Researchers \must report the exact LLM model or tool version, configuration, and experiment date in the \paper. When using quantized models, researchers \should report the quantization level and method. For fine-tuned models, they \must describe the fine-tuning goal, dataset, and procedure in the \paper. Researchers \should include default parameters, explain model choices, compare base- and fine-tuned model using suitable metrics and benchmarks, and share fine-tuning data and weights as \supplementarymaterial (or alternatively justify in the \paper why they cannot share them). 
\end{framed}

\input{_guidelines/02_report-model-version-configuration-and-customizations.tex}

\subsection{Report Tool Architecture Beyond Models (G3)}
\label{sec:report-tool-architecture-beyond-models}

\vspace{0.5\baselineskip}
\begin{framed}
\tldr: \input{_tldr/03_report-tool-architecture-beyond-models-tldr.tex} 
\end{framed}

\input{_guidelines/03_report-tool-architecture-beyond-models.tex}

\subsection{Report Prompts, their Development, and Interaction Logs (G4)}
\label{sec:report-prompts-their-development-and-interaction-logs}

\vspace{0.5\baselineskip}
\begin{framed}
\tldr: \input{_tldr/04_report-prompts-their-development-and-interaction-logs-tldr.tex}
\end{framed}

\input{_guidelines/04_report-prompts-their-development-and-interaction-logs.tex}

\subsection{Use Human Validation for LLM Outputs (G5)}
\label{sec:use-human-validation-for-llm-outputs}

\vspace{0.5\baselineskip}
\begin{framed}
\tldr: \input{_tldr/05_use-human-validation-for-llm-outputs-tldr.tex}  
\end{framed}

\guidelinesubsubsection{Recommendations}

When LLMs are used to automate research or software development tasks previously performed by humans, it is essential to assess the LLM's performance. When existing reference datasets or comparison metrics cannot fully capture the target qualities relevant in the study context, researchers \should rely on human judgment to validate the LLM outputs.

Integrating human participants in the study design will require additional considerations, including a recruitment strategy, annotation guidelines, training sessions, or ethical approvals.
Therefore, researchers \should consider human validation early in the study design, not as an afterthought.
Authors \must clearly define the constructs that the human and LLM annotators evaluate~\cite{DBLP:conf/ease/RalphT18}.
When designing custom instruments to assess LLM output (e.g., questionnaires, scales), researchers \should share their instruments in the \supplementarymaterial.

When the judgment is subjective (i.e. depends on the judge's values or implicit or explicit theories of the phenomenon in question), the LLM output \should be evaluated against judgments aggregated from multiple human judges. In these cases, researchers \should clearly describe their aggregation method and document their reasoning. \added{The preferred method is to randomly order the objects requiring judgment and put them in groups small enough to judge in 1--3 hours. Two or three human experts review the first group together, simultaneously judging, discussing, and creating a set of decision rules documenting their reasoning. Then, the experts iterate among rounds of:}
\begin{itemize}
    \item independently rating one group of objects
    \item calculating inter-rater agreement (IRA) or reliability (IRR)
    \item meeting to discuss disagreements and reach consensus
    \item updating the decision rules so the disagreement will not recur.
\end{itemize}

\added{The goal of this process is to reach some target IRA or IRR (e.g. Krippendorff's $\alpha>0.8$), which indicates that the decision rules are sufficient. If the number of judgments required is large, it is permissible to continue with a single rater per judgment once the IRA/IRR target is reached and sustained for 2--3 groups.} 

Researchers \should provide measures of IRA or IRR, preferably broken down by round. Confounding factors \should be controlled for (e.g., by categorizing participants according to their level of experience or expertise).
Where applicable, researchers \should perform a power analysis to estimate the required sample size, ensuring sufficient statistical power in their experimental design.

Researchers \should use established reference models to compare humans with LLMs.
For example, \citeauthor{Schneider2025ReferenceModel}~\cite{Schneider2025ReferenceModel} outline design considerations for studies comparing LLMs with humans.
%They outline the relevant groups (researchers, experimenters, subjects), activities (study planning, task assignment, results analysis, interpretation), and artifacts (study design, tasks, results, conclusion).
%Further, they illustrate the relationship between these entities, the order of steps, and which design decisions need to be considered.
If studies aim to automate annotating software artifacts using an LLM, researchers \should follow systematic approaches to decide whether and how human annotators can be replaced (e.g. showing that a jury of three LLMs exhibit model-to-model agreement exceeding some threshold).  
% For example, \citeauthor{DBLP:conf/msr/AhmedDTP25}~\cite{DBLP:conf/msr/AhmedDTP25} suggest a method that involves using a jury of three LLMs with 3 to 4 few-shot examples rated by humans, where the model-to-model agreement on all samples is determined using \emph{Krippendorff's Alpha}.
% If the agreement is high (alpha > 0.5), a human rating can be replaced with an LLM-generated one.
% In cases of low model-to-model agreement (alpha $\le$ 0.5), they then evaluate the prediction confidence of the model, selectively replacing annotations where the model confidence is high ($\ge$ 0.8).

Besides generating and modifying content, agentic software development tools such as Claude Code can autonomously call command-line tools or pull in additional information from MCP servers.
For such tools, it makes sense to separate textual output (e.g., file changes) from output related to tool execution.
When evaluating agentic tools, researchers \should assess the feedback that users provided for proposed changes, report statistics on how frequently they accepted the content, and how they modified it.
This agentic human-in-the-loop interaction approach is essentially a built-in human validation, even though the degrees of freedom are larger than in more traditional experiments that use LLMs directly.

\guidelinesubsubsection{Example(s)}

\citeauthor{DBLP:conf/msr/AhmedDTP25}~\cite{DBLP:conf/msr/AhmedDTP25} evaluated how human annotations can be replaced by LLM generated labels.
In their study, they explicitly report the calculated agreement metrics between the models and between humans. They found that three LLMs (GPT-4, Gemini, and Claude) had better agreement on some tasks than human annotators. Meanwhile, \citeauthor{DBLP:conf/icse/XueCBTH24}~\cite{DBLP:conf/icse/XueCBTH24} conducted a controlled experiment to evaluate the impact of ChatGPT on the performance and perceptions of students in an introductory programming course.
They used multiple measures to judge LLM impacts from the human point of view including recording students' screens, evaluating their answers for given tasks, and distributing an opinion survey.

%\citeauthor{hymel2025analysisllmsvshuman}~\cite{hymel2025analysisllmsvshuman} evaluated the capability of ChatGPT-4.0 to generate requirements documents. 
%Specifically, they evaluated two requirements documents based on the same business use case, one document generated with the LLM and one document created by a human expert.
%The documents were then reviewed by experts and judged in terms of alignment with the original business use case (scale of 1-10), requirements completion (Not Complete, Fairly Complete, Fully Complete), and whether they believed it was created by a human or an LLM.
%Finally, they analyzed the influence of the participants' familiarity with AI tools on the study results.
%Participants self-reported their AI tool proficiency (Novice, Intermediate, Advanced, Expert) and usage frequency (Daily, Weekly, Monthly, Never), which were then correlated with their ratings of the requirements documents.

\guidelinesubsubsection{Benefits}

Validating LLMs against human judgments builds confidence in the LLM's accuracy and validity, increasing the trustworthiness of findings, especially when IRA or IRR metrics are reported~\cite{khraisha2024canlargelanguagemodelshumans}. Furthermore, incorporating feedback from human judges may help to improve the LLM or LLM-based tool based on the reported experiences.


%For example, LLM results vary significantly in natural language processing tasks, casting doubt on their reliability and the possibiltiy of replacing human judges~\cite{DBLP:journals/corr/abs-2406-18403}. Moreover, LLMs do not match human annotation in labeling tasks for natural language inference, position detection, semantic change, and hate speech detection~\cite{DBLP:conf/chi/Wang0RMM24}.



\guidelinesubsubsection{Challenges}

Assessing a construct like LLM performance is challenging because ensuring that a construct is defined well and operationalized using an appropriate measurement model requires a deep understanding of (1) the construct, (2) construct validity in general, and (3) instrumentation~\cite{DBLP:journals/tse/SjobergB23,ralph2024teaching}. Comparing an LLM to human judges is typically slower and more expensive than machine-generated measures. More fundamentally, neither human judgment nor machine-generated measures provides an objective ground truth against which LLM accuracy can be firmly determined.   

Human judgments exhibit variability due to differences in experience, expertise, interpretations, and personal biases~\cite{DBLP:journals/pacmhci/McDonaldSF19}. When diverse humans---given clear decision rules corresponding to our construct definitions---rate items reliably, \textit{we simply assume that their reliability implies measurement validity}. It does not. We assume that reliability implies validity because measuring reliability is \textit{much} easier than measuring validity. Indeed, much of the time the best researchers can do is make a conceptual argument for why these judges armed with these decision rules based on this construct \textit{should} produce valid ratings. 


%For example, \citeauthor{hicks_lee_foster-marks_2025} found that racialized software developers \enquote{rated the quality of the output of AI-assisted coding tools significantly lower than other groups}~\cite{hicks_lee_foster-marks_2025}.
% Such biases can be mitigated with a careful selection and combination of multiple judges, but cannot be completely controlled for.

Researchers must weigh the cost and time need for human validation against the expected corresponding improvement in validity over machine-generated measures. 

\guidelinesubsubsection{Study Types}

For \llmusage, researchers \should carefully reflect on the validity of their evaluation criteria.
When an extensively-validated and widely-used benchmark is available (see \benchmarkingtasks), there is little need for human validation.
Where criteria are less clear (e.g., \annotators, \subjects, or \synthesis), or researchers are creating and validating a new benchmark, human validation is important, and this guideline should be followed.
When using \judges, it is common to start with humans who co-create initial rating criteria. 
In developing \newtools, human validation of the LLM output is critical when the output is supposed to match human expectations.

\guidelinesubsubsection{Advice for Reviewers}

\added{Human validation may be the most difficult of our guidelines for reviewers to assess because it will often involve cutting through fallacious arguments. Authors have strong incentives to avoid human validation. However, if LLM output is validated only by comparison with the same or other LLMs, reviewers should demand \textit{quantitative empirical evidence} that such comparison is reliable \textit{and} valid. Similarly, reviewers should demand \textit{quantitative empirical evidence} that any employed benchmarks are reliable \textit{and} valid. Without such evidence, insist on human validation. Showing that several LLMs exhibit high agreement is not sufficient because reliability is necessary but insufficient for validity.}

\added{Next, reviewers will encounter dubious approaches to human validation, such as having a single judge (typically an author) or recruiting dozens of anonymous, untrained, inexperienced gig-workers and aggregating their results on a majority-vote basis. Reviewers should only accept a single judge when the judgments depend only on widely-accepted theories and entails limited value conflict (e.g. tagging method names that contain abbreviations or acronyms). Unfortunately, researchers tend to underestimate the level how theory- and value-laden phenomena are. Any assessment of code quality or tagging of bugs, for example, would be highly dependent on the judge's values and implicit theories~\cite{ralph2024teaching}. }

\added{For multiple judges, reviewers should expect research designs that include techniques known to improve IRA/IRR such as recruiting raters with appropriate education and experience, organizing the rating task into rounds, having raters conduct the initial round of judgments together, having consensus meetings (not voting) after each round in which decision rules are updated. Reviewers should not accept low IRR or IRA (e.g. Krippendorff's $\alpha<0.8$) for research designs absent any of these techniques. Conversely, if authors have embraced best practices for improving IRR/IRA and still obtained middling results (e.g. $0.66<\alpha<0.8$) reviewers should merely insist that low reliability be noted as a limitation.} 

\added{Moving beyond reliability, reviewers should expect authors to \textit{explain} (conceptually, not qualitatively) why human judgments should be valid by comparing construct definitions, directions, decision rules, and judge expertise.}

\added{As with previous guidelines, when important information is missing, reviewers should simply request it be added unless so much is missing that the reviewer cannot assess methodological rigor. For example, failing to provide judge instructions, instruments, decision rules, or construct definitions may prevent assessment of rigor and validity, while missing information about recruitment is less serious. Reviewers should be patient when requesting clearer definitions---the same construct definitions that seem clear to authors often seem opaque to reviewers, and papers should not be rejected over routine clarification requests.}

\subsection{Use an Open LLM as a Baseline (G6)}
\label{sec:use-an-open-llm-as-a-baseline}

\vspace{0.5\baselineskip}
\begin{framed}
\tldr: When researchers can control which LLM a study uses, they \should include an open-weight LLM, ideally an open-source LLM, as a baseline and report inter-model agreement. Researchers \should ensure the open-weight baseline is independently reproducible from their \supplementarymaterial.
 
\end{framed}

\guidelinesubsubsection{Rationale}

Reproducibility depends on access to the model under study.
When research relies exclusively on proprietary models, other researchers cannot independently verify or build upon the findings.
Including an open LLM as a baseline ensures that at least part of the study can be fully replicated.

\guidelinesubsubsection{Recommendations}

Empirical studies using LLMs in SE, especially those that target commercial tools or models, \should incorporate an open LLM as a baseline and report established metrics for inter-model agreement (see \benchmarksmetrics).
We acknowledge that including an open LLM baseline might not always be possible, for example, if the study involves human participants, and letting them work on the tasks using two different models might not be feasible.
Using an open model as a baseline is also not necessary if the use of the LLM is tangential to the study goal.

Open models allow other researchers to verify research results and build upon them, even without access to commercial models.
A comparison of commercial and open models also allows researchers to contextualize model performance.
Researchers \should ensure the open-LLM baseline is independently reproducible from their \supplementarymaterial.

Open LLMs are available from hubs such as \href{https://huggingface.co/}{\emph{Hugging Face}}. They can be self-hosted with frameworks such as \href{https://ollama.com/}{\emph{Ollama}} or \href{https://lmstudio.ai/}{\emph{LM Studio}}, accessed through cloud services such as \href{https://together.ai/}{\emph{Together AI}}, AWS, Azure, and Google Cloud, or routed through aggregators such as \href{https://openrouter.ai/}{\emph{OpenRouter}} that expose many providers behind a single API. For agentic setups, open-source tools such as \href{https://www.continue.dev/}{\emph{Continue}}, \href{https://cline.bot/}{\emph{Cline}}, and \href{https://opencode.ai/}{\emph{opencode}}~\cite{opencode} publish their full agent code, system prompts, and tool catalog. By contrast, vendor-hosted services such as \emph{GitHub Copilot} and \emph{Claude Code} expose only parts of their tooling (e.g., editor extensions, hook examples) and keep their deployed agent loops, system prompts, and tool catalogs proprietary.
% GROUNDING opencode: "OpenCode is an open source AI coding agent."

The term ``open'' can have different meanings in the context of LLMs.
\citeauthor{widder2024open} discuss three types of openness (i.e., transparency, reusability, and extensibility) and what openness can and cannot provide~\cite{widder2024open}.
% GROUNDING widder2024open: "We highlight three main affordances of 'open' AI, namely transparency, reusability, and extensibility"
The \emph{Open Source Initiative} (OSI)~\cite{OSIAI2024} defines open-source AI as having access to everything needed to understand, modify, share, retrain, and recreate the model.
% GROUNDING OSIAI2024: "Sufficiently detailed information about the data used to train the system so that a skilled person can build a substantially equivalent system."

\guidelinesubsubsection{Examples}

An increasing number of studies have adopted open LLMs as baseline models.
For example, \citeauthor{DBLP:journals/corr/abs-2406-09834} evaluated seven advanced LLMs, six of which were open-source, testing 145 API mappings drawn from eight popular Python libraries across 28,125 completion prompts aimed at detecting deprecated API usage in code completion~\cite{DBLP:journals/corr/abs-2406-09834}.
% GROUNDING DBLP:journals/corr/abs-2406-09834: "This study involved seven advanced LLMs, 145 API mappings from eight popular Python libraries, and 28,125 completion prompts."
\citeauthor{DBLP:journals/corr/abs-2408-04430} compared four LLMs on a cross-language code clone detection task~\cite{DBLP:journals/corr/abs-2408-04430}.
% GROUNDING DBLP:journals/corr/abs-2408-04430: "We evaluated four LLMs and one embedding model for cross-lingual code clone detection."
Three evaluated models were open-source.
\citeauthor{DBLP:conf/eicc/GoncalvesSC0MPS25} fine-tuned the open LLM \emph{LLaMA} 3.2 on a refined version of the \emph{DiverseVul} dataset to benchmark vulnerability detection performance~\cite{DBLP:conf/eicc/GoncalvesSC0MPS25}.
% GROUNDING DBLP:conf/eicc/GoncalvesSC0MPS25: "we present a refined version of DiverseVul dataset, which is used to fine-tune a large language model, LLaMA 3.2, for vulnerability detection"
\citeauthor{DBLP:journals/corr/abs-2601-11895} evaluated nine code completion models on the \emph{DevBench} benchmark and included three open-weight models (\emph{DeepSeek-V3}, \emph{DeepSeek-V3.1}, and \emph{Ministral-3B}) alongside commercial frontier models, releasing benchmark, evaluation scripts, and per-model raw completions under an MIT license~\cite{DBLP:journals/corr/abs-2601-11895}.
% GROUNDING DBLP:journals/corr/abs-2601-11895: "Claude 4 Sonnet, Claude 3.7 Sonnet, GPT-4.1 mini, GPT-4.1, GPT-4o, DeepSeek-V3, DeepSeek-V3.1, GPT-4.1 nano, Ministral 3B"
\emph{CodeBERT}, a bimodal transformer pre-trained by Microsoft Research, is published with model weights, source code, and data processing scripts on GitHub~\cite{codebert}. It has been used as an open baseline across diverse SE tasks, including exploit code generation~\cite{DBLP:journals/jss/YangZCZHC23}, vulnerability detection~\cite{DBLP:conf/gaiis/XiaSD24}, code clone detection~\cite{DBLP:conf/kbse/SonnekalbGBM22}, and programming assistance for exception handling~\cite{DBLP:conf/icse/CaiYMMN24}.
% GROUNDING codebert: "This repo provides the code for reproducing the experiments in CodeBERT: A Pre-Trained Model for Programming and Natural Languages."
% GROUNDING DBLP:journals/jss/YangZCZHC23: "we propose a novel template-augmented exploit code generation approach ExploitGen based on CodeBERT"
% GROUNDING DBLP:conf/gaiis/XiaSD24: "this study introduces a software defect detection system leveraging CodeBERT and Bi-LSTM"
% GROUNDING DBLP:conf/kbse/SonnekalbGBM22: "Transformer networks such as CodeBERT already achieve outstanding results for code clone detection in benchmark datasets"
% GROUNDING DBLP:conf/icse/CaiYMMN24: "we design Neurex via a multi-tasking model with the fine-tuning of the large language model CodeBERT for the three above exception handling recommending tasks"

\guidelinesubsubsection{Benefits}

A true open LLM baseline improves reproducibility by exposing model architectures, parameter settings, and ideally training data, enabling independent verification of results.
Such baselines also let researchers compare novel methods against a stable reference point, since proprietary models can silently change between tests.
They also allow inspection of training data (when released) and model behavior, helping identify biases and limitations.
Unlike closed-source alternatives, which can be withdrawn or silently updated, open LLMs remain available for future studies.
They typically avoid the per-use API fees that can constrain research groups with limited budgets.

\guidelinesubsubsection{Challenges}

Open-source LLMs face several challenges:
\begin{itemize}
    \item \emph{Definitional inconsistency.} Many models release only trained weights without training data or methodological details (``open weight'' openness)~\cite{Gibney2024}, which is why we reference the OSI definition in our recommendations~\cite{OSIAI2024}.
    % GROUNDING Gibney2024: "Technology giants such as Meta and Microsoft are describing their artificial intelligence (AI) models as 'open source' while failing to disclose important information about the underlying technology"
    % GROUNDING OSIAI2024: "The complete source code used to train and run the system"
    \item \emph{Performance gap.} Open models often lag behind the most advanced proprietary ``frontier'' models in common benchmarks, making it difficult to demonstrate clear improvements when evaluating new methods using open LLMs alone.
    \item \emph{Hardware demands.} Deploying and experimenting with these models typically requires substantial hardware resources, in particular high-performance GPUs that may be beyond reach for many academic groups.
    \item \emph{Operational complexity.} Unlike APIs provided by proprietary vendors (e.g., the OpenAI API), installing, configuring, and fine-tuning open-source models can be technically demanding.
\end{itemize}

\guidelinesubsubsection{Study Types}

This guideline applies primarily to study types in which the researcher controls which LLM is used.
For \benchmarkingtasks, an open LLM \should be one of the models under evaluation, so that the reported scores can be independently re-run. Where this is not feasible, researchers \should justify the omission and, per \design, ensure the evaluation harness can be used with open models.
For controlled experiments under \llmusage, an open LLM \should be one of the models under test; if the experimental design requires capabilities that only a specific commercial model exhibits, researchers \should acknowledge this as a limitation.
When evaluating \newtools, researchers \should use an open LLM as a baseline whenever it is technically feasible; if integration proves too complex, they \should report the initial benchmarking results of open models.
For \annotators and \judges, researchers \should compare annotation or judgment quality from open vs.\ commercial models to assess the extent to which results depend on a specific proprietary model.
For \synthesis, researchers \should compare synthesis results from open and commercial models to evaluate robustness of the findings.
For \llmusage, using an open LLM as a baseline is often not feasible when the study observes participants using specific commercial tools; in such cases, investigators \should explicitly acknowledge its absence and discuss how this limitation might affect their conclusions.
For \subjects, using an open LLM as a baseline may similarly be impractical if the study design requires the capabilities of a specific model; researchers \should acknowledge this limitation when applicable.

\guidelinesubsubsection{Advice for Reviewers}

Reviewers should distinguish between LLM use that is central to the research (e.g., building LLM-driven SE tools, using LLMs for synthesis) and use that is tangential (e.g., generating recruiting materials). An open LLM baseline is expected only when LLM use is central; otherwise, its absence need not be justified.
When an open baseline is expected, reviewers should look for either the use of an open LLM or a convincing argument for why it is impractical. If authors claim openness, some justification for that characterization is appropriate. Where practical, reviewers should examine whether the replication package contains sufficiently detailed instructions.
Performance differences between open and proprietary models are beyond the study authors' control. Reviewers should focus on whether the open baseline serves the study's methodological purpose rather than on its absolute performance.

\guidelinesubsubsection{See Also}
\begin{itemize}[label=$\rightarrow$]
    \item \seemodelversion: Authors must report version, configuration, and parameters for open models, not just commercial ones.
    \item \seedesign: Self-hosting an open model adds hardware, hosting, and infrastructure to document.
    \item \seebenchmarksmetrics: Inter-model agreement metrics make open-vs-commercial comparisons interpretable.
    \item \seelimitationsmitigations: When using an open LLM as a baseline is impractical, authors must report this as a limitation.
\end{itemize}


\subsection{Use Suitable Baselines, Benchmarks, and Metrics (G7)}
\label{sec:use-suitable-baselines-benchmarks-and-metrics}

\vspace{0.5\baselineskip}
\begin{framed}
\tldr: \input{_tldr/07_use-baselines-benchmarks-and-metrics-tldr.tex}  
\end{framed}

\guidelinesubsubsection{Recommendations}

Benchmarks, baselines, and metrics play an important role in assessing the effectiveness of LLMs and LLM-based tools. \added{A \textit{benchmark} is a ``standard tool for the competitive evaluation and comparison of competing systems or components according to specific characteristics, such as performance, dependability, or security''~\cite{Kistowski2015Benchmark} }
Benchmarks can be used to assess LLM performance on specific tasks such as code summarization, code generation, \added{and static analysis}.
Meanwhile, ``a \textit{metric} is a method, algorithm, or procedure for assigning one or more numbers to a phenomenon''~\cite{ralph2024teaching}. 
A baseline represents a reference point for the measured LLM.
Since LLMs require substantial hardware resources, baselines allow research to compare LLMs against traditional algorithms with lower computational costs.

When selecting benchmarks and metrics, it is important to fully understand the benchmark tasks, what exactly is being measured, and how it relates to the (often latent) variables researchers actually care about. If one or more metrics or benchmarks are used, researchers:
\begin{itemize}
    \item \must briefly justify (in the \paper) why selected benchmarks and metrics are suitable for the given task or study. 
    % \item \must discuss the reliability and validity (especially construct validity) of selected metrics and benchmarks
    \item \should summarize the structure and tasks of the selected benchmark(s), including the programming language(s) and descriptive statistics such as the number of contained tasks and test cases;
    \item \should discuss the limitations of the selected benchmark(s) (e.g. widely used benchmarks such as \emph{HumanEval} \cite{DBLP:journals/corr/abs-2107-03374} and \emph{MBPP} \cite{DBLP:journals/corr/abs-2108-07732} Python functions assess a very specific part of software development, which is not representative of the full breadth of SE work~\cite{Chandra2025benchmarks}).
    \item \added{\should} include an example of a task and the corresponding test case(s) to illustrate the structure of the benchmark.
\end{itemize}

If multiple benchmarks exist for the same task, researchers \should compare performance between benchmarks.
When selecting only a subset of all available benchmarks, researchers \should use the most specific benchmarks given the context.

Furthermore, researchers \should check whether a less resource-hungry approach (e.g., for static analysis tasks or program repair) can serve as a baseline. If so, the LLM or LLM-based tool should be compared to such baselines using suitable metrics. Even if LLM-based tools outperform baselines, researchers \should discuss whether the resources consumed justify the (often marginal) improvements~\cite{DBLP:journals/cacm/Menzies25}.


To compare traditional and LLM-based approaches or different LLM-based tools, researchers \should report established metrics whenever possible, as this allows secondary research.
They can, of course, report additional metrics that they consider appropriate.

In the following, we briefly discuss two common metrics used for generation tasks: \emph{BLEU-N} and \emph{pass@k}.
\emph{BLEU-N} \cite{DBLP:conf/acl/PapineniRWZ02} is a similarity score based on n-gram precision between two strings, ranging from $0$ to $1$.
Values close to $0$ represent dissimilar content, values closer to $1$ represent similar content, indicating that a model is more capable of generating the expected output.
%The score measures how many n-grams from the generated output of a model appear in the target string.
%In isolation, this approach favors shorter over longer generated sequences because it is more likely to have fewer n-grams in an extensive target sequence than more.
%To mitigate this, the so-called brevity penalty is introduced to disincentivize short sequences when having longer target sequences.
\emph{BLEU-N} has multiple variations.
\emph{CodeBLEU} \cite{DBLP:journals/corr/abs-2009-10297} and \emph{CrystalBLEU} \cite{DBLP:conf/kbse/EghbaliP22} are the most notable ones tailored to code.
They introduce additional heuristics such as AST matching.
As mentioned, researchers \must motivate why they chose a certain metric or variant thereof for their particular study.

The metric \emph{pass@k} reports the likelihood that a model correctly completes a code snippet at least once within \emph{k} attempts.
To our knowledge, the basic concept of \emph{pass@k} was first reported by \citeauthor{DBLP:journals/corr/abs-1906-04908} to evaluate code synthesis~\cite{DBLP:journals/corr/abs-1906-04908}.
They named the concept \emph{success rate at B}, where B denotes the trial ``budget.''
The term \emph{pass@k} was later popularized by \citeauthor{DBLP:journals/corr/abs-2107-03374} as a metric for code generation correctness~\cite{DBLP:journals/corr/abs-2107-03374}.
The exact definition of correctness varies depending on the task.
For code generation, correctness is often defined based on test cases: A passing test implies that the solution is correct.
The resulting pass rates range from $0$ to $1$.
A pass rate of $0$ indicates that the model was unable to generate a single correct solution within $k$ tries; a rate of $1$ indicates that the model successfully generated at least one correct solution in $k$ tries.

The \emph{pass@k} metric is defined as:

\begin{equation}
\text{pass@}k = 1 - \frac{\binom{n-c}{k}}{\binom{n}{k}}
\end{equation}

where $n$ is the total number of generated samples per prompt, $c$ is the number of correct samples among  $n$, and $k$ is the number of samples.

Choosing an appropriate value for \emph{k} depends on the downstream task of the model and how end-users interact with it.
% k= 1 -- pass@1
A high pass rate for \emph{pass@1} is highly desirable in tasks where the system presents only one solution or if a single solution requires high computational effort.
For example, code completion depends on a single prediction, since end users typically see only a single suggestion.
%In automated code refactoring, a large context might lead to high computational costs, making the calculation of a single solution expensive.
%Thus, having a high pass rate at pass@1 is crucial for such tasks.
% k = >=2 -- pass@{>=2}
Pass rates for higher $k$ values (e.g., $2$, $5$, $10$) indicate how well a model can solve a given task in multiple attempts.
%For downstream tasks that allow multiple solutions, e.g. via user interaction, strong performance at $k > 1$ can be justified. 
%For instance, a user selecting the correct test case from multiple suggestions allows for some model errors.
%In code search, finding the most appropriate good snippet for a given context, a reranking algorithm mitigate low pass rates at low k-values.
The \emph{pass@k} metric is frequently referenced in papers on code generation \cite{DBLP:journals/corr/abs-2308-12950, DBLP:journals/corr/abs-2401-14196, DBLP:journals/corr/abs-2409-12186, DBLP:journals/corr/abs-2305-06161}.
However, \emph{pass@k} is not a universal metric suitable for all generation tasks.
It requires a binary notion of correctness, making the metric appropriate for code synthesis evaluated via unit tests, but unsuitable for open-ended generation tasks such as comment generation, where multiple valid outputs exist.

If a study evaluates an LLM-based tool for supporting humans, a relevant metric is the acceptance rate, meaning the ratio of all accepted artifacts (e.g., test cases, code snippets) in relation to all artifacts that were generated and presented to the user.
Another way of evaluating LLM-based tools is calculating inter-model agreement (see also Section \href{/guidelines/#use-an-open-llm-as-a-baseline}{Use an Open LLM as a Baseline}).
This allows researchers to assess how dependent a tool's performance is on specific models and versions.
Metrics used to measure inter-model agreements include general agreement (percentage), \emph{Cohen's kappa}, and \emph{Scott's Pi coefficient}.

Due to LLM non-determinism, researchers \should repeat experiments to  assess (statistically) model or tool performance, e.g., using the arithmetic mean, median, confidence intervals, standard deviations, or more advanced statistical approaches~\cite{DBLP:conf/nips/AgarwalSCCB21, bjarnason2026randomnessagenticevals}.

From a measurement perspective, researchers \should reflect on the theories, values, and measurement models on which the benchmarks and metrics they have selected for their study are based.
For example, the phenomena labeled ``bugs'' in a large open dataset of software bugs is according to a certain theory of what constitutes a bug, and from the values and perspectives of the people who labeled the dataset. Reflecting on the context in which these labels were assigned and discussing whether and how the labels generalize to a new study context is crucial.

\guidelinesubsubsection{Example(s)}

According to \citeauthor{10.1145/3695988}, the main problem types for LLMs are classification, recommendation, and generation, each of which requires different metrics. In their systematic review,
\citeauthor{hu2025assessingadvancingbenchmarksevaluating} \cite{hu2025assessingadvancingbenchmarksevaluating} categorized 191 LLM benchmarks according to the specific SE task they address, making the paper a valuable resource for identifying existing metrics. 
Metrics used to assess generation tasks include \emph{BLEU}, \emph{pass@k}, \emph{Accuracy}, \emph{Accuracy@k}, and \emph{Exact Match}.
The most common metric for recommendation tasks is \emph{Mean Reciprocal Rank}.
For classification tasks, classical machine learning metrics such as \emph{Precision}, \emph{Recall}, \emph{F1-score}, and \emph{Accuracy} are often reported.

%\guidelinesubsubsection{Exact Match}
%\emph{Exact match} is another metric that calculates the percentage of replicas a model can produce.
%If the model is able to produce the exact target sequence, a score of 1 is awarded; otherwise 0.
%Compared to BLEU-N, \emph{exact match} is a stricter measurement.
%Due to its strictness, reported scores are usually low.

%he choice of metrics strongly depends on the specific task.
%Consequently, papers that address different research questions employ different metrics.
%For example, \citeauthor{DBLP:conf/nips/LiuXW023} assesses LLM models for their correctness in code generation using \emph{pass@k}~\cite{DBLP:conf/nips/LiuXW023}.
%However, \emph{pass@k} is not an universal metric suitable for every problem at hand.
%For instance, for creative or open-ended tasks such as comment generation, pass@k is not a suitable metric since not only one single correct result exists.
%Thus, \citeauthor{DBLP:conf/icse/GengWD00JML24} evaluates LLMs for code comment generation using BLEU, ROUGE-L and METEOR as evaluation %metrics~\cite{DBLP:conf/icse/GengWD00JML24}.
%\citeauthor{DBLP:journals/ase/LiuJZNLL25} investigates the performance of LLMs on automated refactorings tasks and proposes a novel metric tailored to the unique characteristics of this problem \cite{DBLP:journals/ase/LiuJZNLL25}.

%Code snippets provided by benchmarks contain untrusted code, making sandboxing recommendable.
%To avoid potential security risks, execute untrusted code in an isolated environment, such as a virtual machine or containerization, instead of your own environment.

Benchmarks used for code generation incldue \emph{HumanEval} (available on \href{https://github.com/openai/human-eval}{GitHub}) \cite{DBLP:journals/corr/abs-2107-03374}, \emph{MBPP} (available on \href{https://huggingface.co/datasets/google-research-datasets/mbpp}{Hugging Face}) \cite{DBLP:journals/corr/abs-2108-07732},
% Both benchmarks consist of code snippets written in Python sourced from publicly available repositories.
% Each snippet consists of four parts: a prompt containing a function definition and a corresponding description of what the function should accomplish, a canonical solution, an entry point for execution, and test cases.
% The input of the LLM is the entire prompt.
% The output of the LLM is then evaluated against the canonical solution using metrics or against a test suite.
\emph{ClassEval} (available on \href{https://github.com/FudanSELab/ClassEval}{GitHub}) \cite{DBLP:journals/corr/abs-2308-01861}, \emph{LiveCodeBench} (available on \href{https://github.com/LiveCodeBench/LiveCodeBench}{GitHub}) \cite{DBLP:journals/corr/abs-2403-07974}, and \emph{SWE-bench} (available on \href{https://github.com/swe-bench/SWE-bench}{GitHub}) \cite{DBLP:conf/iclr/JimenezYWYPPN24}.
An example of a code translation benchmark is \emph{TransCoder}~\cite{DBLP:journals/corr/abs-2006-03511} (available on \href{https://github.com/facebookresearch/CodeGen}{GitHub}). 

% Lukas Boehme: Add examples for other categories: classification and recommendation
Given the generative nature of LLMs, most benchmarks focus on generation tasks, but benchmarks for classification and recommendation tasks also exist.
For classification tasks such as vulnerability detection, \emph{DiverseVul} (available on \href{https://github.com/wagner-group/diversevul}{GitHub})~\cite{DBLP:conf/raid/0001DACW23} provides a dataset of vulnerable and non-vulnerable functions.
They evaluate the binary classification using  \emph{Accuracy}, \emph{Precision}, \emph{Recall}, \emph{F1-score} and \emph{False Positive Rate}.
For recommendation tasks such as code search, \emph{CodeSearchNet} (available on \href{https://github.com/github/CodeSearchNet}{GitHub})~\cite{DBLP:journals/corr/abs-1909-09436} contains code-documentation pairs across multiple programming languages.
The recommendations are evaluated using \emph{Mean Reciprocal Rank}.

\guidelinesubsubsection{Benefits}

Benchmarks and metrics improve reproducibility and comparability across studies, leading to faster improvements in the cumulative body of knowledge. It is simply easier to assess a new system against one or a few benchmarks than hundreds of competing systems, and when most systems are assessed against the same benchmarks and metrics, we can see which approaches work best (at least, on those benchmarks). This comparison enables progress tracking; for example, when researchers iteratively improve a new LLM-based tool and test it against benchmarks after significant changes. For practitioners, leaderboards (published benchmark results of models) support selecting models for downstream tasks. \added{Meanwhile, baselines help us understand just how much LLMs improve performance over non-LLM-enabled alternatives.} 

\guidelinesubsubsection{Challenges}

\added{Computer scientists have a long history of inventing oversimplified, unvalidated metrics (e.g. lines of code, CPU time) and simply claiming that they are reasonable proxies for complicated, value laden, under-theorized, often multidimensional, latent variables (e.g. system size, environmental sustainability). In scientific fields that take measurement more seriously, researchers expect that metrics and instruments are backed up by foundational theories (e.g. we accept telescopes based on our understanding of the physics of light) or empirical evidence (e.g. we accept questionnaire scales for the Big Five personality traits based on extensive quantitative and qualitative investigations of their psychometric properties). The fundamental challenge with AI metrics and benchmarks is that many of them have neither firmly understood theoretical underpinnings nor extensive empirical validation of their construct, measurement, and ecological validity.} 

For example, metrics such as \emph{pass@k} do not reflect qualitative aspects of the generated source code, including its maintainability or readability.
However, these aspects are critical for many downstream tasks.
Moreover, benchmarks are isolated test sets and may not fully represent real-world applications.
For example, benchmarks such as \emph{HumanEval} synthesize code based on written specifications.
However, such explicit descriptions are rare in real-world applications.
Thus, evaluating model performance with benchmarks might not reflect real-world tasks and end-user performance.

Benchmarks and metrics also have generalizability problems. Prominent LLM benchmarks such as \emph{HumanEval} and \emph{MBPP} use Python, so researchers can optimize for Python's idiosyncrasies, creating illusory improvements that do not generalize to other languages. Similarly, the metric \emph{BLEU-N} is a syntactic metric, so code can score highly without being executable. The metric \emph{Exact Match}, meanwhile, does not account for functional equivalence of syntactically different code. Both \emph{BLEU-N} and \emph{Exact Match} are influenced by code formatting, which confounds their intended use. 

Execution-based metrics such as \emph{pass@k} directly evaluate correctness by running test cases, but they require a setup with an execution environment.
When researchers observe unexpected values for certain metrics, the specific results should be investigated in more detail to uncover potential problems.


% Many closed-source models, such as those released by \emph{OpenAI}, achieve exceptional performance on certain tasks but lack transparency and reproducibility~\cite{DBLP:conf/nips/00110ZZDJLHL24, DBLP:journals/corr/abs-2308-01861, DBLP:journals/corr/abs-2406-15877}.
% Benchmark leaderboards, particularly for code generation, are led by closed-source models~\cite{DBLP:journals/corr/abs-2308-01861, DBLP:journals/corr/abs-2406-15877}.
% While researchers \should compare performance against these models, they must consider that providers might discontinue them or apply undisclosed pre- or post-processing beyond the researchers' control (see \openllm).
%The lack of transparency introduces challenges in analyzing mistakes, limiting the understanding of model failures.
%Open-source models offer greater transparency since the entire process is under the researcher's control; however, they require appropriate hardware.

\added{Finally, benchmark data contamination---where the benchmark itself is part of the training data---may lead to artificially high performance if the model remembers the solution from the training data rather than actually solving the task based on the input data~\cite{DBLP:journals/corr/abs-2406-04244}. Therefore, for proprietary LLMs that do not release their training data, researchers \should consider using human validation (Section~\ref{sec:use-human-validation-for-llm-outputs}), curating new data, or refactoring existing data, or~\cite{DBLP:journals/corr/abs-2406-04244}.} 

\added{Creating and curating new uncontaminated benchmark datasets produces unseen benchmark data.
This can be achieved by collecting data after a specified cutoff date, ensuring that models trained before this date could not have been exposed to the benchmark. To ensure transparency, researchers \must disclose the data sources and their collection dates for each benchmark release. However, this temporal approach may be impractical because it requires continuous updates as new models might include the benchmark in future training datasets. Keeping the benchmark private prevents inclusions in training sets, yet it requires trust in the benchmark creator and a system to execute the benchmark without leaking the benchmark.}

\added{Refactoring existing data restructures benchmarks by generating new prompts or augmenting samples across multiple dimensions. Thus, pure memorization by the model is ineffective. However, this only partially addresses contamination since the model may have learned underlying concepts from other versions of the benchmark.}


%\removed{Models may become overfitted to the training data, leading to poor performance on unseen data. When optimizing for a metric, the underlying model might adjust its parameters to improve the metric, thereby failing to generalize. This can result in a model performing exceptionally well on a benchmark but struggling to hold up the performance in unforeseen settings, such as real-world scenarios.}

\guidelinesubsubsection{Study Types}

This guideline \must be followed for all study types that automatically evaluate the performance of LLMs or LLM-based tools.
The design of a benchmark and the selection of appropriate metrics are highly dependent on the specific study type and research goal.
Recommending specific metrics for specific study types is beyond the scope of these guidelines, but \citeauthor{hu2025assessingadvancingbenchmarksevaluating} provide a good overview of existing metrics for evaluating LLMs~\cite{hu2025assessingadvancingbenchmarksevaluating}.
In addition, our Section \humanvalidation provides an overview of the integration of human evaluation in LLM-based studies. 
For \annotators, the research goal might be to assess which model comes close to a ground truth dataset created by human annotators.
Especially for open annotation tasks, selecting suitable metrics to compare LLM-generated and human-generated labels is important.
In general, annotation tasks can vary significantly.
Are multiple labels allowed for the same sequence? Are the available labels predefined, or should the LLM generate a set of labels independently?
Due to this task dependence, researchers \must justify their metric choice, explaining what aspects of the task it captures together with known limitations.
If researchers assess a well-established task such as code generation, they \should report standard metrics such as \emph{pass@k} and compare the performance between models.
If non-standard metrics are used, researchers \must state their reasoning.

\guidelinesubsubsection{Advice for Reviewers}

\added{Reviewers should expect manuscripts to: 
\begin{enumerate}
    \item clearly identify the constructs or variables that the study ought to measure in principle (e.g. LLM performance, quality of generated code), including independent variable(s), dependent variable(s) and (both physical and statistical) control variables;
    \item present their measurement model; that is, identify exactly which metrics, benchmarks, or baselines are used and how they relate to constructs or variables the study aims to measure;
    \item explicitly justify the selection of any metrics, benchmarks, and baselines used
    item contemplate \textit{in detail} the assumptions, reliability, and validity---\textit{especially construct validity} of each benchmark and metric;
    \item clearly articulate any limitations regard construct and measurement validity (see also Section~\ref{sec:report-limitations-and-mitigations}).
\end{enumerate}}

\added{As in previous guidelines, if some specific information about a baseline, etc. is missing, reviewers should simply ask authors to add it. In contrast, reviewers must not tolerate vague descriptions that conflate what researchers and practitioners actually care about (e.g. effectiveness, efficiency, quality, etc.) with specific methods of counting something ostensibly indicative of these broader concerns. Reviewers should be wary of popularity arguments; while the whole point of benchmarks is that standardizing evaluation is beneficial, ubiquity does not imply validity or appropriateness for any given context. Similarly, reviewers must not tolerate ambivalence toward construct validity or obfuscation of validity concerns.  In summary, for this guideline, reviewers should look for manuscripts to convey a solid understanding of construct and measurement validity by explaining and justifying selected measurement models. }

\subsection{Report Limitations and Mitigations (G8)}
\label{sec:report-limitations-and-mitigations}

\vspace{0.5\baselineskip}
\begin{framed}
\tldr: \input{_tldr/08_report-limitations-and-mitigations-tldr.tex}   
\end{framed}

\guidelinesubsubsection{Rationale}

When using LLMs for empirical studies in SE, researchers face unique challenges and potential limitations that can influence the validity, reliability, and reproducibility of their findings~\cite{DBLP:conf/icse/SallouDP24}.
% GROUNDING DBLP:conf/icse/SallouDP24: "researchers must still be careful as numerous intricate factors can influence the outcomes of experiments involving LLMs"
Researchers must openly discuss these limitations and explain how their impact was mitigated.
These limitations are relative to current LLM capabilities and tool architectures; speculating about future improvements is beyond the scope of a paper's limitation section.
Nevertheless, risk management and threat mitigation \should be planned during study design, not as an afterthought.

\guidelinesubsubsection{Recommendations}
Researchers \must clearly present the limitations of their work without defensiveness or obfuscation. These limitations may concern external, internal, and construct validity; reliability and reproducibility; and ethical, regulatory, and environmental concerns.

We follow the standard SE convention of organizing limitations by external, internal, and construct validity, plus reliability~\cite{DBLP:journals/ese/RunesonH09,ralph2021empiricalstandardssoftwareengineering}, extended below with ethical, regulatory, and environmental concerns specific to LLM research. For studies that adopt qualitative analytic methods (the use of LLMs in reflexive qualitative analysis is itself contested, see \annotators), researchers \should use the trustworthiness criteria of credibility, transferability, dependability, and confirmability instead~\cite{guba1981criteria}. When deterministic reproducibility is structurally unattainable (e.g., SaaS-based models with opaque versioning), researchers \should adopt the same trustworthiness criteria to substantiate dependability and confirmability of findings.
% GROUNDING DBLP:journals/ese/RunesonH09: "This scheme distinguishes between four aspects of the validity, which can be summarized as follows: Construct validity ... Internal validity ... External validity ... Reliability"
% GROUNDING ralph2021empiricalstandardssoftwareengineering: "General Quality Criteria (e.g. internal validity, construct validity, credibility, resonance)"
% GROUNDING guba1981criteria: "The four terms credibility, transferability, dependability and confirmability were established as the naturalist's equivalents for the more conventional rigor criteria of internal validity, external validity, reliablity and objectivity."

\paragraph{External Validity.}
The primary threats to external validity are:
\begin{itemize}
    \item \emph{Cross-model transfer limitations}: Results obtained with one LLM or family of LLMs may not generalize to others due to differences in training data, architecture, and post-training procedures (e.g., fine-tuning and reinforcement learning from human feedback); see \openllm for using an open LLM as a comparison baseline.
    \item \emph{Configuration sensitivity}: Results may not generalize beyond the specific configuration(s) tested (e.g., decoding parameters, system prompt, or context settings); see \design for an overview.
    \item \emph{Tool-architecture specificity}: Tools built around vendor-specific APIs or features (e.g., function calling, structured output, or context-window size) may not transfer to other models without substantial re-engineering (see \design).
    \item \emph{Limited domain coverage}: Studies often focus on a narrow set of programming languages, task types, or application domains, limiting generalizability to other SE contexts.
    \item \emph{Limited participant and rater diversity}: Study participants (e.g., developers) and human raters validating LLM outputs may not represent the broader population in terms of expertise, geographic location, or cultural background (see \humanvalidation for guidance on rater diversity in value-laden constructs).
    \item \emph{Research-to-practice gap}: Developers using the same tools outside researcher-supervised conditions may obtain results that differ from those reported in the study.
    \item \emph{Cross-time instability}: Performance of proprietary models can change over time, leading to non-generalizable study outcomes~\cite{DBLP:journals/corr/abs-2307-09009, doi:10.1148/radiol.232411} (see \modelversion for version and fingerprint reporting that enables tracking such drift).
    % GROUNDING DBLP:journals/corr/abs-2307-09009: "the behavior of the "same" LLM service can change substantially in a relatively short amount of time"
    % GROUNDING doi:10.1148/radiol.232411: "Our results suggest performance drift between the March and June snapshots of GPT-3.5 and GPT-4."
\end{itemize}

Generalizability is particularly critical for proprietary and non-deterministic systems whose behavior is subject to drift (i.e., silent changes in model output over time). Researchers \must discuss the limitations and mitigations of external validity.
Mitigations include \emph{triangulation} across multiple models (e.g., proprietary and open), independent datasets, and complementary metrics; \emph{sensitivity analysis} that varies LLM configurations, prompts, architecture decisions, datasets, and where applicable participant backgrounds; and performing \emph{longitudinal re-runs} (see \emph{Reliability \& Reproducibility} below), which also helps detect cross-time instability.

\paragraph{Internal Validity.}
The primary threats to internal validity are:
\begin{itemize}
	\item \emph{Data leakage and contamination}: Inter-dataset duplication can produce training-evaluation overlap, yielding overly optimistic results (see \benchmarksmetrics).
	\item \emph{Evaluation data entering model-improvement pipelines}: Evaluation samples can unintentionally feed retraining or fine-tuning, especially in longitudinal studies involving LLMs.
	\item \emph{Incomplete architecture, prompt, or pipeline reporting}: Undisclosed components introduce hidden confounders (see \design).
\end{itemize}

As transparency on training data is limited for LLMs, researchers \must discuss potential data leakage effects and their impact on results. Concrete mitigations for contamination (e.g., post-cutoff benchmark construction, held-out subsets, canary strings) are discussed in \benchmarksmetrics.

\paragraph{Construct Validity.}
The primary threats to construct validity are:
\begin{itemize}
	\item \emph{Metric-construct mismatch}: Traditional metrics such as BLEU or ROUGE may miss SE-specific aspects such as functional correctness or behavioral equivalence (see \benchmarksmetrics).
	\item \emph{Construct under-specification}: If a construct lacks an operational definition, neither automated metrics nor human raters can apply it consistently.
	\item \emph{Reliability without validity}: High inter-rater or inter-model agreement does not imply that the measurement captures the intended construct; a reliable LLM can be reliably inaccurate (see \humanvalidation).
	\item \emph{Benchmark scope limitations}: Benchmarks commonly ignore runtime behaviors, security implications, readability, testability, and maintainability, yielding results that may not transfer to realistic development settings.
	\item \emph{Capability confounding}: Benchmark performance can blend the target capability with unrelated capabilities such as output format compliance, instruction following, or prompt format sensitivity, inflating apparent scores (see \benchmarksmetrics).
	\item \emph{Over-reliance on benchmark-specific metrics}: Optimizing for a single benchmark may produce dataset-specific shortcuts that pass the benchmark without exhibiting the capability it was designed to test, overstating real-world utility.
	\item \emph{Prompt sensitivity}: Small changes in instructions, formatting, or in-context examples can substantially shift what the LLM appears to measure, making the operationalized construct unstable across prompt variants (see \design).
	\item \emph{Judge biases}: LLM and human judges exhibit systematic biases such as position bias, verbosity bias, and preferences for richly formatted or citation-rich outputs regardless of correctness (see \humanvalidation).
\end{itemize}

If constructs are based on subjective interpretations, purely automated metrics are insufficient. Researchers \must discuss how they ensured quality of subjective results, similarly to qualitative research. The primary mitigation is \emph{human validation} of subjective constructs following quality criteria known from qualitative research (see \humanvalidation).

\paragraph{Reliability \& Reproducibility.}
The primary threats to reliability and reproducibility are:
\begin{itemize}
    \item \emph{Non-deterministic outputs}: Identical prompts and configurations can yield different outputs across runs due to factors such as floating-point arithmetic, batching, and stochastic decoding strategies.
    \item \emph{Infrastructure dependence}: Results may vary depending on the hardware, software stack, and hosting environment used; vendor-imposed quotas, throttling, or pricing changes can further prevent re-running experiments at the original scale, making exact replication challenging across different infrastructure setups.
    \item \emph{Resource inequality}: LLM research is resource-intensive and remains predominantly in the domain of private companies or well-funded research institutions~\cite{schwartz2020green, ahmed2023industry}, excluding researchers from under-resourced institutions.
    % GROUNDING schwartz2020green: "the financial cost of the computations can make it difficult for academics, students, and researchers, in particular those from emerging economies, to engage in deep learning research"
    % GROUNDING ahmed2023industry: "industry increasingly dominates the three key ingredients of modern AI research: computing power, large datasets, and highly skilled researchers"
\end{itemize}

Researchers \must discuss the measures taken to increase reliability and reproducibility.
However, non-deterministic reproducibility is not inherently disqualifying. The trustworthiness criteria introduced above apply particularly to SaaS-based LLM research, where providers frequently deprecate model versions without guaranteeing stable behavior.
Mitigations include providing \emph{replication packages} that cover prompt and architecture specifications, model outputs, and representative examples for partial replicability (ideally accompanied by an implementation using an open model for long-term stability), and performing \emph{longitudinal re-runs} with statistical analyses. When deterministic reproduction is structurally impossible, researchers \should consider \emph{methodological trustworthiness measures} such as triangulation, reflexivity, audit trails, and peer debriefing as complementary measures.

\paragraph{Ethical \& Regulatory Boundaries.}
The primary concerns for ethical and regulatory matters are:
\begin{itemize}
    \item \emph{Use of sensitive or proprietary data}: Studies involving proprietary code, confidential business data, or personally identifiable information may face restrictions on data sharing that limit reproducibility.
    \item \emph{Jurisdictional obligations}: Data protection regulations such as GDPR or CCPA and institutional policies may impose constraints on data collection, processing, and sharing in LLM-based studies.
    \item \emph{Implicit model bias}: Especially for qualitative research, LLMs might \enq{reinforce dominant paradigms and biases} and \enq{identify, replicate and reinforce dominant language and patterns}~\cite{jowsey2025reject} (see \humanvalidation).
    % GROUNDING jowsey2025reject: "Failure to recognize these limitations of GenAI risks analyses that reinforce dominant paradigms and biases."
\end{itemize}

Studies involving sensitive data \must discuss data governance mechanisms tailored toward LLM environments, compliant with applicable jurisdictional obligations. If applicable, researchers \should discuss how model biases potentially impact the study outcomes and how those biases were evaluated.

\paragraph{Environmental \& Sustainability Constraints.}
The primary environmental and sustainability concerns are:
\begin{itemize}
    \item \emph{Energy consumption}: With growing model size, the environmental impact of experiments with LLMs increases, and the substantial energy costs of LLM experiments warrant consideration in study design~\cite{DBLP:conf/acl/StrubellGM19}.
    % GROUNDING DBLP:conf/acl/StrubellGM19: "These accuracy improvements depend on the availability of exceptionally large computational resources that necessitate similarly substantial energy consumption."
    \item \emph{Trade-off between repetition and sustainability}: Repeating experiments increases reliability but also energy consumption, requiring trade-offs during study design.
\end{itemize}

Researchers \should justify the LLM's resource consumption against the benefits over traditional approaches.
Mitigations include \emph{energy preservation} and \emph{cost accounting}. \emph{Energy preservation} involves selecting smaller or newer, less resource-intensive models and applying techniques such as input/output token reduction, model pruning, quantization, or knowledge distillation~\cite{mitu2024hidden} where feasible. Carbon footprint estimation is desirable, but still difficult. \emph{Cost accounting} tracks resource consumption by reporting tokens, service costs, or hardware specifications.
% GROUNDING mitu2024hidden: "Techniques such as model pruning, quantisation, and knowledge distillation can reduce the computational resources needed for training and inference, thereby decreasing the overall carbon footprint"

\guidelinesubsubsection{Examples}

\citeauthor{DBLP:conf/icse/SallouDP24} catalog three categories of LLM-specific threats to validity (i.e., closed-source models, implicit data leakage, and reproducibility) and pair each with concrete mitigation strategies (e.g., versioned model archives, metamorphic test data, multiple replication runs with variability metrics, and detailed execution metadata)~\cite{DBLP:conf/icse/SallouDP24}.
% GROUNDING DBLP:conf/icse/SallouDP24: "potential threats to the validity of LLM-based research including issues such as closed-source models, possible data leakage between LLM training data and research evaluation, and the reproducibility of LLM-based findings"
\citeauthor{DBLP:conf/icse/Du0WWL0FS0L24} pair each threat in their \emph{ClassEval} evaluation with a concrete mitigation: manually constructing the benchmark with multiple annotators to limit data leakage, piloting prompts on held-out tasks to control for prompt sensitivity, and reporting greedy-decoding results to control for non-determinism~\cite{DBLP:conf/icse/Du0WWL0FS0L24}.
% GROUNDING DBLP:conf/icse/Du0WWL0FS0L24: "One potential threat is the data leakage between our benchmark and model training data, thus we manually construct the benchmark ClassEval."

\guidelinesubsubsection{Benefits}

Transparent reporting of limitations and mitigations helps readers calibrate confidence in the findings, makes explicit which threats were addressed and which remain open, and documents design decisions that other authors can borrow or refine. It also keeps a paper's claims proportionate to its evidence.

\guidelinesubsubsection{Challenges}

Identifying limitations one is not already aware of is the hardest part of writing a threats section, particularly for methodological threats outside the team's primary expertise. Publication and reviewing norms can pressure authors to downplay weaknesses, while page limits make exhaustive treatment impractical. Threat lists that recite generic LLM-research issues (e.g., model bias, non-determinism, or contamination) without showing how each one applies to specific design choices in this study leave reviewers unable to tell which risks actually applied.

The threats to validity framework itself is contested within SE. \citeauthor{DBLP:journals/infsof/VerdecchiaELRS23} argue that threats sections too often read as \enq{laundry-lists}~\cite{DBLP:journals/infsof/VerdecchiaELRS23}, \citeauthor{DBLP:conf/esem/LagoRSV24} corroborated this empirically across a decade of ICSE Distinguished Paper Award winners~\cite{DBLP:conf/esem/LagoRSV24}, and \citeauthor{DBLP:journals/tosem/RobillardAEGLNNSSS24} argue for refocusing the discussion on study design trade-offs rather than the standard validity categories~\cite{DBLP:journals/tosem/RobillardAEGLNNSSS24}.
% GROUNDING DBLP:journals/infsof/VerdecchiaELRS23: "TTV sections seem to be mostly formulated as “laundry-lists” of potential threats, lacking a thorough contextualization in the specifics of the study at hand"
% GROUNDING DBLP:conf/esem/LagoRSV24: "the selected papers, despite being selected as best papers in the ICSE conference, reported dominantly shallow or none TTV analysis"
% GROUNDING DBLP:journals/tosem/RobillardAEGLNNSSS24: "First, threats to validity sections are often boilerplate text produced by rote."

\guidelinesubsubsection{Study Types}

Researchers \must follow this guideline for all study types. Transparently reporting limitations and mitigations is a universal requirement, but specific concerns vary by study type.
For \annotators, researchers \must discuss potential biases in label assignment, label reliability limitations, and sensitivity of annotations to prompt wording and model choice.
For \judges, researchers \must address measurement validity concerns, known biases such as position bias or verbosity bias, and the extent to which LLM judgments align with human expert assessments.
For \synthesis, researchers \must discuss the risk of contextual misinterpretation, potential loss of nuance in summarized or aggregated outputs, and reflexivity limitations inherent in using an LLM for qualitative interpretation.
For \subjects, researchers \must discuss the fundamental inability of LLMs to truly simulate human behavior, the risk of stereotype amplification, and the limited ecological validity of simulated responses.
For \llmusage, researchers \must discuss generalizability constraints across different tools and user populations, and acknowledge how observed usage patterns may not transfer to other contexts.
For \newtools, researchers \must discuss replicability constraints arising from dependencies on commercial models, the impact of model updates on tool behavior, and limitations of the evaluation setup.
For \benchmarkingtasks, researchers \must discuss potential data contamination, benchmark scope limitations, capability confounding, and the extent to which benchmark performance generalizes to real-world tasks.

\guidelinesubsubsection{Advice for Reviewers}

Reviewers should verify that the limitation section is comprehensive and appropriate for the specific study type, checking that:
(1) limitations address the specific threats relevant to the study type (e.g., label reliability for annotation studies, simulation fidelity for studies using LLMs as subjects);
(2) mitigations are concrete and correspond to identified limitations rather than being generic statements;
(3) the impact of LLM non-determinism on findings is discussed;
(4) generalizability constraints, across models, configurations, time periods, and populations, are acknowledged.
When important limitations are missing, reviewers should request they be added. The absence of a limitation section, or one that is formulaic or insufficiently specific, is a more serious concern than any individual missing limitation and may warrant a major revision.

\guidelinesubsubsection{See Also}
\begin{itemize}[label=$\rightarrow$]
    \item \seemodelversion: Authors can re-run with different models or configurations to check whether the results depend on those specific choices.
    \item \seedesign: Triangulation across architectures and prompts is one mitigation strategy.
    \item \seetraces: Stored session traces serve as a baseline against which authors can monitor LLM behavior drift over time.
    \item \seebenchmarksmetrics: Benchmark and metric choices are one source of construct-validity threats authors must discuss.
    \item \seeopenllm: An open LLM baseline mitigates cross-model transfer concerns by providing an independently reproducible comparison.
    \item \seehumanvalidation: When automated metrics cannot validly measure a construct, human validation is an alternative.
\end{itemize}

\begin{table}[tb]
\caption{Rationale and key recommendations for each guideline. \iconM~=~\must, \iconS~=~\should.}
\label{tab:rationale-recommendations}
\small
\begin{tabularx}{\textwidth}{@{}l>{\raggedright\arraybackslash}X>{\raggedright\arraybackslash}X@{}}
\toprule
\textbf{ID} & \textbf{Rationale} & \textbf{Core Recommendations} \\
\midrule
G1 & Transparency enables informed assessment of scope and limitations. &
\iconM Declare which LLM, how it was used, and where in the research process. \\
\addlinespace
G2 & Reproducibility requires precise identification of the system used in a study. &
\iconM Report exact version, date, configuration, and fine-tuning details. \\
\addlinespace
G3 & Complex systems require holistic documentation beyond the model itself. &
\iconM Describe full architecture, agent roles, and infrastructure. \\
\addlinespace
G4 & Prompts shape LLM behavior and must be verifiable. &
\iconM Publish all prompts, context files, and development strategy.
\newline \iconS Share interaction logs where feasible. \\
\addlinespace
G5 & Automated metrics alone cannot ensure validity of subjective constructs. &
\iconS Validate against human judgment with inter-rater reliability, where applicable. \\
\addlinespace
G6 & Reproducibility depends on access to the model under study. &
\iconS Include open LLM baseline; provide replication package. \\
\addlinespace
G7 & Meaningful evaluation requires reasoned valid measurement. &
\iconM Justify metric and benchmark choices; discuss their validity.
\newline \iconS Repeat experiments and report descriptive statistics. \\
\addlinespace
G8 & Honest acknowledgment of threats strengthens a study. &
\iconM Report threats by validity type.
\newline \iconS Employ mitigations where possible. \\
\bottomrule
\end{tabularx}
\end{table}


\section{Conclusion}

LLM non-determinism, opaque training data, and rapidly evolving models
threaten the reproducibility and replicability of empirical SE studies.
This paper presents a taxonomy of seven study types that organizes how
LLMs are used in SE research, together with eight guidelines for
designing and reporting such studies.
Each guideline distinguishes requirements (\must) from recommended
practices (\should) and is contextualized by the study types it
applies to.
Our guidelines recommend that researchers:
\begin{enumerate*}[label=(\arabic*)]
\item declare LLM usage and role;
\item report model versions, configurations, and customizations;
\item document the tool architecture beyond the model;
\item disclose prompts, their development, and, where feasible,
interaction logs;
\item validate LLM outputs with humans;
\item include an open LLM as a baseline;
\item use suitable baselines, benchmarks, and metrics; and
\item articulate limitations and mitigations, including costs, potential
biases, and environmental impact.
\end{enumerate*}
An applicability matrix maps guidelines to study types, and a reporting
checklist (Section~\ref{sec:checklist}) supports authors and reviewers.
The study types and guidelines were derived collaboratively by
22~co-authors, following a process similar to that of guidelines in
other disciplines (see, e.g.,~\cite{Begg1996}).

Following the quality dimensions of the \emph{SIGSOFT Empirical
Standards}~\cite{ralph2021empiricalstandardssoftwareengineering}, we
consider our guidelines
(1)~\emph{comprehensive}, as they cover seven study types that we
identified and were derived from 22 co-authors' collective expertise;
(2)~\emph{useful}, as they provide actionable recommendations with an
applicability matrix and a concrete reporting checklist
(Section~\ref{sec:checklist});
(3)~\emph{well-argued}, as each guideline is motivated by a dedicated
rationale and supported by references to the literature; and
(4)~\emph{integrated with previous work}, as they build on and
complement established empirical SE standards while addressing
LLM-specific challenges.
Adopting these practices will not eliminate LLM non-determinism, but it
should make claims auditable and results easier to replicate and compare
across studies.
Because models and tools evolve, we maintain the guidelines as a living
resource and invite the community to refine and extend them at
\href{https://llm-guidelines.org/}{llm-guidelines.org}.

\section*{Acknowledgments}

Our study types and guidelines are based on discussion sessions with researchers at the \emph{2024 International Software Engineering Research Network} (ISERN) meeting and at the \emph{2nd Copenhagen Symposium on Human-Centered Software Engineering AI} (supported by the Carlsberg Foundation with grant CF24-0693 and the Alfred P. Sloan Foundation with grant G-2024-22586 to Daniel Russo).
LLM-based tools, including ChatGPT (GPT-5) and Claude Code (Opus 4.6), were used interactively for language editing and structural revision suggestions across all sections of this paper. All LLM-generated suggestions were reviewed and revised by the authors.

\begin{appendices}

\section{Reporting Checklist}
\label{sec:checklist}

% NOTE: Each \item below carries a non-rendered comment citing the source
% sentence(s) in the corresponding _guidelines/0X_*.tex file. Update both
% the item and the comment when guideline text changes.
% Ordering rule (per CLAUDE.md): within each section, sort by severity
% (\iconM before \iconS), then by general before \condition{}-tagged
% (conditional) items, then by reporting location
% (\paper -> unspecified -> \supplementarymaterial).
% Selection rule: every \must and \mustnot statement in the guideline bodies
% appears in (or is absorbed into) an item. A \should gets its own item when
% it names a standalone reporting action or artifact; a \should that refines,
% conditions, or illustrates an included item is folded into that item (quote
% it in the item's comment), and interpretation cautions or reflective-practice
% advice without a discrete reportable artifact stays out. Document deliberate
% exclusions in a comment (see the power-analysis note under Human Validation).

The following checklist, inspired by CONSORT~\citep{Schulz2010}, summarizes actionable items from the guidelines.
% GROUNDING Schulz2010: "The CONSORT 2010 Statement is this paper including the 25 item checklist in the table (Table 1) and the flow diagram (Figure 1)."
The checklist is organized along typical paper sections.
Items marked \iconM are requirements (\must), and items marked \iconS are recommendations (\should).
Each item references its source guideline by short name.
Items annotated with \paper or \supplementarymaterial indicate where we expect the information to be reported.
Unmarked items may be reported either in the paper or as supplementary material.
Items prefixed with a bracketed tag apply only to studies with that characteristic (e.g., \condition{fine-tuning}, \condition{agents}).\ifpaper\else{} Readers can hover each tag to read its description and use the filter panel to hide items that do not apply to their study.\fi
Beyond these characteristic-tagged items, each guideline's \emph{Study Types} subsection lists study-type-specific recommendations where applicable.

\paragraph{\textbf{Introduction}}
\begin{itemize}[label={},leftmargin=*]
    \item \iconM Disclose any use of LLMs in the empirical study, specifying which LLM, how, and where it was used~(\refdeclareusage).
    % Declare Usage Recommendations: "researchers \must clearly declare that an LLM was used"; specifies which/how/where in the research process. Folds the multiple-roles refinement: "each role \should be declared separately".
    \item \iconS Report in the \paper the purpose of using LLMs, the tasks they automate, and the expected benefits~(\refdeclareusage).
    % Declare Usage Recommendations: "researchers \should report the exact purpose of using an LLM in a study, the tasks it was used to automate, and the expected benefits in the \paper".
\end{itemize}

\paragraph{\textbf{Research Design and Methods}}\mbox{}\\
\noindent\textit{Model Selection and Configuration}
\begin{itemize}[label={},leftmargin=*]
    \item \iconM Report in the \paper the exact LLM model or tool version, the configuration, and the date of study execution~(\refmodelversion).
    % Model Version Recommendations: "Researchers \must document in the \paper which model or tool version they used... along with the date of study execution and the parameters they configured that affect output generation".
    \item \iconM \condition{fine-tuning} For fine-tuned models, describe in the \paper the fine-tuning goal, the dataset, and the procedure~(\refmodelversion).
    % Model Version Recommendations: "If a model was fine-tuned, researchers \must describe the fine-tuning goal..., the fine-tuning procedure..., and the fine-tuning dataset... in the \paper".
    \item \iconS Report default parameters and explain model and version choices~(\refmodelversion).
    % Model Version Recommendations: "researchers \should report all configuration values, even if they used the defaults"; "Researchers \should motivate in the \paper why they selected certain models, versions, and configurations". Folds the temperature-0 sub-case: "this motivation \should be stated explicitly"; the \shouldnot caution against treating temperature 0 as a reproducibility guarantee is an interpretation caution, not an item.
    \item \iconS Report checksums and additional model properties where available; for commercial tools, openly acknowledge their reproducibility limits~(\refmodelversion).
    % Model Version Recommendations: "Checksums and fingerprints \should be reported"; "other properties such as the context window size... \should be reported"; for commercial tools, "researchers \should report what is available and openly acknowledge limitations that hinder reproducibility". Wording compresses "checksums and fingerprints" to "checksums" for symmetry with the summary-table row; the guideline body retains the full "checksums and fingerprints" phrasing.
    \item \iconS \condition{quantization} For quantized models, report the quantization level (e.g., 4-bit, 8-bit) and method (e.g., GPTQ or AWQ)~(\refmodelversion).
    % Model Version Recommendations: "When using quantized models, researchers \should report the quantization level... and method".
    \item \iconS \condition{fine-tuning} Compare base and fine-tuned models using suitable metrics and benchmarks; share fine-tuning data and weights as \supplementarymaterial (or justify in the \paper why they cannot be shared)~(\refmodelversion).
    % Model Version Recommendations: "Suitable benchmarks and metrics \should be used to compare the base model with the fine-tuned model"; "Researchers \should either share the fine-tuning dataset as part of the \supplementarymaterial or explain in the \paper why the data cannot be shared. The same applies to the fine-tuned model weights".
    \item \iconS \condition{commercial-models} Include an open LLM as a baseline when using commercial models and report inter-model agreement~(\refopenllm).
    % Open LLM Recommendations: "Empirical studies using LLMs in SE, especially those that target commercial tools or models, \should incorporate an open LLM as a baseline and report established metrics for inter-model agreement". Placed in Model Selection and Configuration because the trigger is a model-choice condition (use of commercial models), even though the recommendation itself involves a baseline.
\end{itemize}

\noindent\textit{System and Prompt Design}
% Within each severity: general items first (architecture, prompting strategy,
% prompt reuse, standalone setup, prompt-publication; then S meta, hosting, and
% prompt-development items, then open-source release), then \condition{}-tagged
% items in topic order (prompting, prompt sharing, configuration, tools and
% agents, retrieval, benchmarking, timing).
\begin{itemize}[label={},leftmargin=*]
    \item \iconM Describe in the \paper the full architecture of LLM-based tools, including the role of the LLM, interactions with other components, and overall system behavior~(\refdesign).
    % Design Recommendations: "Researchers \must clearly describe the tool architecture and what exactly the LLM... contributes"; for the topic-specific paragraphs that follow, "the \paper \must contain a high-level description of any reported component, with full details provided as \supplementarymaterial".
    \item \iconM Specify whether zero-shot, one-shot, or few-shot prompting was used~(\refdesign).
    % Design Prompting Strategy: "Researchers \must specify whether zero-shot, one-shot, or few-shot prompting was used" (no reporting location in the body).
    \item \iconM Specify prompt reuse across models and configurations~(\refdesign).
    % Design Prompt Development: "researchers \must make clear which prompts were used for which models in which versions and with which configuration".
    \item \iconM If the LLM was used in a standalone setup, with prompts sent directly to a model and no pre- or post-processing, state this explicitly~(\refdesign).
    % Design Pipelines and Complex Systems: "If the LLM is used in a \emph{standalone setup}, with prompts sent directly to a model via an API and no pre-processing of inputs or post-processing of outputs, researchers \must state this explicitly".
    \item \iconM Publish all prompts or, when using templates, prompt templates with representative instances, including their structure, content, formatting, and variable components, as \supplementarymaterial; include representative examples in the \paper~(\refdesign).
    % Design Prompt Reporting / Supplementary: "researchers \must report all prompts used in an empirical study"; "When prompt templates are used, researchers \must report them alongside representative instances"; "The complete set \must be made publicly available as \supplementarymaterial, with representative examples in the \paper itself".
    \item \iconM \condition{few-shot} For few-shot prompts, explain in the \paper how the examples were selected~(\refdesign).
    % Design Prompting Strategy: "For few-shot prompts, researchers \must explain in the \paper how the examples were selected".
    \item \iconM \condition{dynamic-prompts} For dynamically generated prompts, document the code or rules that assemble each prompt from runtime inputs~(\refdesign).
    % Design Prompt Reporting: "When prompts are generated dynamically, researchers \must document the code or rules that assemble each prompt from runtime inputs".
    \item \iconM \condition{restricted-sharing} For partially shared prompts, anonymize personal identifiers, replace proprietary code with placeholders, and clearly highlight modified sections~(\refdesign).
    % Design Prompt Reporting: "For prompts that can be partially shared, researchers \must anonymize personal identifiers, replace proprietary code with placeholders, and clearly highlight modified sections".
    \item \iconM \condition{context-files} Describe in the \paper any configuration mechanisms used (e.g., context files such as \texttt{CLAUDE.md} or \texttt{AGENTS.md}, skills, subagents, hooks, settings, rules)~(\refdesign).
    % Design Context Files and Agent Configuration: "researchers \must describe which configuration mechanisms were used".
    \item \iconM \condition{tool-use} Summarize in the \paper which tools were exposed to the model~(\refdesign).
    % Design Tool Catalog and MCP Servers: "Researchers \must summarize in the \paper which tools were exposed to the model".
    \item \iconM \condition{agents} If autonomous agents are used, specify agent roles, reasoning frameworks, and communication flows~(\refdesign).
    % Design Agentic Systems: "researchers \must additionally describe the agents' roles..., whether the system is single-agent or multi-agent, how the agents interact with external tools and users, and the reasoning framework used". Folds the subagent refinement: "researchers \should also describe the delegation pattern".
    \item \iconM \condition{agents} For agentic systems that use external tools, distinguish the model's reasoning and outputs, its tool calls, and its interactions with users, the environment, or other agents~(\refdesign).
    % Design Agentic Systems: "researchers \must distinguish three kinds of activity: (1) the model's reasoning, planning, and outputs; (2) tool calls (e.g., to APIs, databases, file systems, or ... (MCP) servers); (3) interactions with users, the environment, or other agents".
    \item \iconM \condition{context-augmentation} For retrieval-augmented generation (RAG) or related methods, describe how external data was retrieved, stored, and selected for inclusion in the model's context~(\refdesign).
    % Design Pipelines and Complex Systems (RAG sub-case): "researchers \must additionally describe how external data was retrieved, stored ..., and selected for inclusion in the model's context".
    \item \iconM \condition{benchmarking} Describe the evaluation harness and infrastructure when it goes beyond bare model API calls (e.g., custom sandboxing, orchestration layers, or post-processing pipelines)~(\refdesign).
    % Design Study Types (benchmarking): "researchers \must describe the evaluation harness and infrastructure when it goes beyond bare model API calls (e.g., custom sandboxing, orchestration layers, or post-processing pipelines)".
    \item \iconM \condition{benchmarking} For pre-defined prompts from existing benchmarks, specify the benchmark version, and disclose and justify any modifications to the prompts or evaluation setup~(\refdesign).
    % Design Study Types (benchmarking): "Researchers using pre-defined prompts (e.g., \emph{HumanEval}, \emph{SWE-Bench}) \must specify the benchmark version and any modifications made to the prompts or evaluation setup"; "If prompt tuning, RAG, or other methods were used to adapt prompts, researchers \must disclose and justify those changes".
    \item \iconM \condition{latency-sensitive} For time-sensitive measurements, clarify whether local infrastructure or cloud services were used, including the specific hardware for local hosting (e.g., GPU model and VRAM) or the service tier for cloud APIs~(\refdesign).
    % Design Recommendations: "For \emph{time-sensitive measurements}, the hosting choice... can substantially affect results, so researchers \must clarify whether local infrastructure or cloud services were used, including the specific hardware for local hosting (e.g., GPU model and VRAM) or the service tier for cloud APIs".
    \item \iconS Justify substantive architectural choices where alternatives existed (e.g., agentic framework, tool catalog)~(\refdesign).
    % Design Recommendations: "Researchers \should justify substantive architectural choices where alternatives existed (e.g., why a particular agentic framework or tool catalog was selected)".
    \item \iconS Describe how the models were hosted and accessed~(\refdesign).
    % Design Recommendations: "Researchers \should describe how the models were hosted and accessed." The conditional \must for time-sensitive measurements appears as the latency-sensitive item above.
    \item \iconS Describe prompt development rationale and selection process~(\refdesign).
    % Design Prompt Development: "Researchers \should explain in the \paper how they developed the prompts and why they decided to follow certain prompting strategies"; "If multiple versions of a prompt were tested, researchers \should describe how these variations were evaluated".
    \item \iconS Report prompt evolution and any LLM-suggested refinements~(\refdesign).
    % Design Prompt Development: "researchers \should at least summarize the prompt evolution"; "Researchers \should report any instances in which LLMs were used to suggest prompt refinements".
    \item \iconS Where legally possible, release the source code of the implementation under an open-source license~(\refdesign).
    % Design Recommendations: "Where legally possible..., researchers \should release the source code of their implementation under an open-source license".
    \item \iconS \condition{few-shot} Include the concrete few-shot examples in the \supplementarymaterial~(\refdesign).
    % Design Prompting Strategy: "For few-shot prompts, researchers... \should include the concrete examples in the \supplementarymaterial".
    \item \iconS \condition{participant-prompts} For user-authored prompts, describe how they were collected and analyzed~(\refdesign).
    % Design Prompt Reporting: "For studies involving human participants who create or modify prompts, researchers \should describe how these prompts were collected and analyzed".
    \item \iconS \condition{long-prompts} Document input handling and token optimization strategies when prompts are long or complex~(\refdesign).
    % Design Prompting Strategy: "researchers \should describe the strategies they used to handle input length constraints"; "Token optimization measures... \should also be documented".
    \item \iconS \condition{restricted-sharing} If full prompt disclosure is not feasible, provide summaries or examples~(\refdesign).
    % Design Prompt Reporting: "When confidentiality (e.g., industry-partner agreements) prevents full publication, researchers \should publish summaries and representative examples instead".
    \item \iconS \condition{ensemble} For ensemble architectures, explain in the \paper the coordination logic between models~(\refdesign).
    % Design Pipelines and Complex Systems (Ensembles): "researchers \should describe the architecture connecting them: the routing logic that determines which model handles which input, model interactions, and the output combination strategy". The general paper/supp split applies to the topic-specific paragraphs. Coordination-logic specifics are S.
    \item \iconS \condition{context-augmentation} Report data preprocessing, versioning, and update frequency for stored data used for context augmentation~(\refdesign).
    % Design Pipelines and Complex Systems (RAG sub-case): "The data used for retrieval \should also be reported, including its preprocessing, versioning, and update frequency".
    \item \iconS \condition{context-files} Include all configuration artifacts (context files, skill folders, subagent files, hooks, settings, rules) as \supplementarymaterial~(\refdesign).
    % Design Context Files and Agent Configuration: researchers "\should publish all configuration artifacts as \supplementarymaterial".
    \item \iconS \condition{tool-use} Include the tool catalog (names with purposes), tool schemas, and connected MCP servers as \supplementarymaterial~(\refdesign).
    % Design Tool Catalog and MCP Servers: "Researchers \should include a complete list with names and purposes, tool schemas, and the names of any connected MCP servers as \supplementarymaterial".
    \item \iconS \condition{benchmarking} Design the evaluation harness so it is usable with open models~(\refdesign).
    % Design Study Types (benchmarking): "the harness \should be usable with open models". Carries the replicability load after \openllm's benchmarking carve-out was downgraded from \must to \should.
\end{itemize}

\noindent\textit{Session Traces}
\begin{itemize}[label={},leftmargin=*]
    \item \iconM Where tool-native trace formats are used (e.g., Claude Code's session transcripts, LangGraph's state logs), describe the file format and report the tool version~(\reftraces).
    % Traces Runtime Traces: "Where tool-native formats are used (e.g., Claude Code's session transcripts, LangGraph's state logs), researchers \must describe the file format and report the tool version".
    \item \iconM \condition{restricted-sharing} For traces containing sensitive information, anonymize personal identifiers, replace proprietary code with placeholders, and clearly highlight modified sections~(\reftraces).
    % Traces Interaction Logs: "For traces containing sensitive information, researchers \must anonymize personal identifiers, replace proprietary code with placeholders, and clearly highlight modified sections".
    \item \iconS Use an open trace format with a documented schema where it fits (e.g., the OpenTelemetry GenAI semantic conventions or OpenInference)~(\reftraces).
    % Traces Runtime Traces: "Researchers \should use an open format with a documented schema. Emerging standards such as the OpenTelemetry GenAI semantic conventions... or OpenInference... are preferred where they fit".
    \item \iconS Include full interaction logs (prompts and responses) as \supplementarymaterial if privacy and confidentiality can be ensured~(\reftraces).
    % Traces Interaction Logs: "Researchers \should report the full interaction logs... as part of their \supplementarymaterial". (For \llmusage specifically the matrix elevates this to \must.)
    \item \iconS \condition{agents} For agentic systems, include interaction logs covering human-in-the-loop exchanges with the agent (feedback, approvals, refinements) as \supplementarymaterial~(\reftraces).
    % Traces Interaction Logs: "For agentic systems, interaction logs cover the human-facing exchanges with the agent, including human-in-the-loop feedback, approval or rejection decisions, and iterative refinements. These \should also be reported as \supplementarymaterial".
    \item \iconS \condition{agents} For agentic systems, report the complete runtime trace as \supplementarymaterial, including for each entry the tool or artifact name, arguments, result, and ordering, and which configured artifacts (skills, context files, subagents) were activated~(\reftraces).
    % Traces Runtime Traces: "Researchers \should report the complete runtime trace as \supplementarymaterial, including for each entry the tool or artifact name, arguments (if any), result, and ordering relative to surrounding interaction-log entries". The merged trace covers both per-call detail and artifact activations.
    \item \iconS \condition{agents} For agentic systems, report any plans the system exposes as \supplementarymaterial~(\reftraces).
    % Traces Agentic Plans: "researchers \should report any plans the system exposes as \supplementarymaterial".
\end{itemize}

\noindent\textit{Benchmarks and Metrics}
% Within each severity: general items first, then \condition{}-tagged items.
% Open LLM cross-link (open LLM baseline) sits at the end of the section so it does
% not interrupt the Benchmarks & Metrics validity arc.
\begin{itemize}[label={},leftmargin=*]
    \item \iconM Justify in the \paper all benchmark and metric choices~(\refbenchmarks).
    % Benchmarks & Metrics Recommendations: researchers "\must briefly justify in the \paper why each selected benchmark and metric is suitable for the given task or study".
    \item \iconM Discuss in the \paper the reliability and validity, especially construct validity, of the selected benchmarks and metrics~(\refbenchmarks).
    % Benchmarks & Metrics Recommendations: researchers "\must discuss the reliability and validity---especially construct validity---of those choices (see \limitationsmitigations)".
    \item \iconM Explain in the \paper why the selected metrics are suitable for the specific study; prior adoption in related work alone is not sufficient justification~(\refbenchmarks).
    % Benchmarks & Metrics Recommendations: "researchers \must motivate why they chose a certain metric or variant thereof for their particular study"; "Prior adoption alone does not constitute sufficient justification; researchers \mustnot justify metric choices solely by citing their use in prior work".
    \item \iconM \condition{latency-sensitive} Report latency when it can affect study outcomes (e.g., interactive user studies, latency comparisons)~(\refbenchmarks).
    % Benchmarks & Metrics Recommendations: "Latency \must be reported when it can affect study outcomes". Study Types: for \llmusage, "in interactive setups where response times can influence participant behavior, researchers \must report observed latency"; for \benchmarkingtasks, "when researchers compare LLMs or tools on latency, they \must report the infrastructure used to produce the measurements".
    \item \iconM \condition{new-benchmark} For new or updated benchmarks, disclose data sources and collection dates for each release~(\refbenchmarks).
    % Benchmarks & Metrics Challenges: "researchers \must disclose data sources and collection dates for each release".
    \item \iconS Provide an operational definition of the phenomenon the benchmark is intended to measure, including its scope and any sub-components~(\refbenchmarks).
    % Benchmarks & Metrics Recommendations (new): researchers "\should provide an operational definition of the phenomenon the benchmark is intended to measure ..., including its scope and any sub-components".
    \item \iconS Summarize benchmark structure, task types, and limitations~(\refbenchmarks).
    % Benchmarks & Metrics Recommendations: researchers "\should summarize the structure and tasks of the selected benchmark(s)" and "\should discuss the limitations of the selected benchmark(s)". Folds the illustration refinement: "include an example of a task and the corresponding test case(s)".
    \item \iconS Identify the capabilities a benchmark conflates with the target phenomenon, isolate the target where possible, and acknowledge remaining confounders as construct-validity threats~(\refbenchmarks).
    % Benchmarks & Metrics Recommendations (new): researchers "\should identify which capabilities a benchmark conflates, isolate the target capability where possible ..., and acknowledge remaining confounders as threats to construct validity".
    \item \iconS Perform an error analysis: categorize the failures observed and report their relative frequency; report failures that cluster on confounding capabilities as construct-validity threats~(\refbenchmarks).
    % Benchmarks & Metrics Recommendations (new): researchers "\should perform an error analysis by categorizing the failures observed and reporting the relative frequency of each category". "If failures concentrate on inputs that demand capabilities other than the target capability ..., this is a construct-validity threat".
    \item \iconS Describe and justify the sampling strategy used to select problems for inclusion in the benchmark~(\refbenchmarks).
    % Benchmarks & Metrics Recommendations (new): researchers "\should describe and justify the sampling strategy used to select problems for inclusion in the benchmark". Conditional follow-up about non-probability sampling appears as a separate \condition{}-tagged bullet below.
    \item \iconS Justify the number of experiment repetitions, for example through a power analysis or by monitoring convergence of descriptive statistics~(\refbenchmarks).
    % Benchmarks & Metrics Recommendations: "researchers \should justify their chosen number of repetitions, for example, through a power analysis... or by monitoring the convergence of descriptive statistics".
    \item \iconS \condition{non-probability-sampling} For non-probability sampling (e.g., convenience), discuss the implications for the generalizability of conclusions~(\refbenchmarks).
    % Benchmarks & Metrics Recommendations (new): "if non-probability sampling (e.g., convenience) is used, discuss its implications for the generalizability of conclusions".
    \item \iconS \condition{new-benchmark} For new or released benchmarks, adopt contamination-prevention mechanisms: held-out subset, canary strings, and pre-exposure investigation against common training corpora~(\refbenchmarks).
    % Benchmarks & Metrics Recommendations (new): researchers creating or releasing a new benchmark "\should adopt concrete contamination-prevention mechanisms: maintaining a held-out subset of items for ongoing, uncontaminated evaluation; embedding canary strings ...; and investigating whether the benchmark's source materials may already appear in common LLM training corpora".
    \item \iconS \condition{multi-rater-scoring} For ratings that vary across raters or runs (human raters, LLM-as-judge), report the distribution of ratings per item rather than only aggregated point estimates~(\refbenchmarks).
    % Benchmarks & Metrics Recommendations (new): "When scoring uses ratings that vary across raters or runs ..., researchers \should report the distribution of ratings per item rather than only aggregated point estimates or exact-match agreement".
\end{itemize}

\noindent\textit{Human Validation}
\begin{itemize}[label={},leftmargin=*]
    \item \iconM \condition{human-validation} If using human validation, define in the \paper the measured construct (e.g., usability, maintainability) and describe the measurement instrument~(\refhumanvalidation).
    % Human Validation Study Design Considerations: "Authors \must clearly define in the \paper the constructs that the human and LLM annotators evaluate".
    \item \iconM \condition{human-validation} When developing or adapting measurement instruments, share them~(\refhumanvalidation).
    % Human Validation Study Design Considerations: "When designing custom instruments to assess LLM output (e.g., questionnaires, scales), researchers \must share these instruments".
    \item \iconM \condition{human-validation} When LLMs replace humans in research tasks, explain whether and how the replacement is justified~(\refhumanvalidation).
    % Human Validation Recommendations (Replacing Human Judgment): "In such cases, researchers \must explain whether and how the replacement is justified". Echoed in the Study Types umbrella sentence for \annotators, \judges, \synthesis, and \subjects.
    \item \iconS Consider human validation early in the study design and build on established reference models for human-LLM comparison~(\refhumanvalidation).
    % Human Validation Recommendations: "researchers \should consider human validation early in the study design, not as an afterthought"; "Researchers \should use established reference models to compare humans with LLMs".
    \item \iconS \condition{human-validation} When LLMs replace humans in research tasks, report the systematic approach used to justify the replacement, including inter-model and model-to-human agreement~(\refhumanvalidation).
    % Human Validation Recommendations (Replacing Human Judgment): "researchers \should follow systematic approaches to support this judgment, e.g., showing that a jury of three LLMs exhibits inter-model agreement... such as Krippendorff's $\alpha>0.8$"; "researchers \should additionally validate a sample against human experts"; "Researchers \should justify why the LLM is methodologically appropriate".
    \item \iconS \condition{human-validation} Validate LLM judgments against human judgment, report aggregation methods, and assess human-LLM agreement~(\refhumanvalidation).
    % Human Validation Subjective Judgment and Agreement: "the same artifacts \should be judged independently by both the LLM and a panel of human experts, and the LLM judgments \should then be compared against an aggregated human reference. Researchers \should clearly describe their aggregation method"; "Researchers \should report measures of IRA or IRR".
    \item \iconS \condition{human-validation} Discuss and, where feasible, control for confounding factors~(\refhumanvalidation).
    % Human Validation Subjective Judgment and Agreement: "Confounding factors \should be discussed and, where feasible, controlled for (e.g., by categorizing participants according to their level of experience or expertise)". Power analysis is intentionally excluded here: Human Validation source uses plain "may" ("Where applicable, researchers may perform a power analysis...") and the next sentence explicitly notes "limited established guidance specific to determining sample sizes for LLM-human comparisons" - so it does not warrant a \should-level checklist item.
    \item \iconS \condition{human-validation}\condition{subjective-constructs} For value-laden or culturally contingent constructs, describe rater demographics beyond expertise and discuss potential demographic biases~(\refhumanvalidation).
    % Human Validation Subjective Judgment and Agreement (new): "For value-laden or culturally contingent constructs ..., researchers \should describe rater demographics beyond expertise (e.g., geographic, linguistic, or professional background) and discuss potential demographic biases in rater recruitment and instructions".
    \item \iconS \condition{agents} When evaluating agentic tools, assess the feedback users provided on the agent's proposed actions, and report how frequently proposals were accepted and how they were modified~(\refhumanvalidation).
    % Human Validation Agentic Tools: "researchers \should assess the feedback that users provided on the agent's proposed actions (e.g., file edits, command executions), report statistics on how frequently they accepted the proposals, and how they modified them".
\end{itemize}

\noindent\textit{Reproducibility, Ethics, and Resources}
\begin{itemize}[label={},leftmargin=*]
    \item \iconM \condition{restricted-sharing} For studies involving sensitive data, discuss data governance mechanisms compliant with applicable jurisdictional obligations~(\reflimitations).
    % Limitations Ethical \& Regulatory: "Studies involving sensitive data \must discuss data governance mechanisms tailored toward LLM environments, compliant with applicable jurisdictional obligations".
    \item \iconS Justify LLM usage in light of its resource demands~(\reflimitations).
    % Limitations Environmental \& Sustainability: "Researchers \should justify the LLM's resource consumption against the benefits over traditional approaches".
    \item \iconS Ensure the open-LLM baseline is independently reproducible from the \supplementarymaterial~(\refopenllm).
    % Open LLM Recommendations: "Researchers \should ensure the open-LLM baseline is independently reproducible from their \supplementarymaterial". This replaces the former full-replication-package item: the broad guidance is an unmarked mitigation in Limitations (covered by the mitigation-strategies item under Limitations and Threats to Validity), while the RFC-marked statement is this open-LLM-scoped \should.
    \item \iconS \condition{restricted-sharing} Where full sharing of prompts, traces, or datasets is not feasible, share representative examples for partial replicability~(\reflimitations).
    % Limitations Reliability \& Reproducibility Mitigations: "Replication Packages: cover prompt and architecture specifications, model outputs, and representative examples for partial replicability, accompanied by an implementation using an open model for long-term stability". The mitigation bullet has no RFC 2119 marker; \iconS is interpretive, treating it as \should-level. Concrete instances of "if full sharing not feasible, share representative examples" appear with explicit \should in Model Version (fine-tuning data), Design (prompts), and Traces (interaction logs); this item is the cross-artifact summary.
\end{itemize}

\paragraph{\textbf{Results}}
\begin{itemize}[label={},leftmargin=*]
    \item \iconS Repeat experiments due to the inherent non-determinism of LLMs and report the result distribution using descriptive statistics~(\refbenchmarks).
    % Benchmarks & Metrics Recommendations: "Due to LLM non-determinism, researchers \should repeat experiments and report descriptive statistics of model or tool performance".
    \item \iconS Use traditional (non-LLM) baselines for comparison where possible~(\refbenchmarks).
    % Benchmarks & Metrics Recommendations: "researchers \should check whether a less resource-intensive approach... can serve as a baseline. If so, the LLM or LLM-based tool \should be compared with such baselines using suitable metrics".
    \item \iconS Report established metrics to make study results comparable; additional metrics may be reported where appropriate~(\refbenchmarks).
    % Benchmarks & Metrics Recommendations: "researchers \should report established metrics whenever possible, as this allows secondary research. They can report additional metrics".
    \item \iconS \condition{comparing-models} If comparing models or tools, use appropriate inferential statistics (e.g., hypothesis tests, effect sizes) rather than relying solely on summary statistics~(\refbenchmarks).
    % Benchmarks & Metrics Recommendations: "If comparing models or tools, researchers \should use appropriate inferential statistics... rather than relying solely on differences in means or other summary statistics". Downgraded from \must to \should: see CHANGELOG entry on the disclosure-vs-methodology design pattern.
\end{itemize}

\paragraph{\textbf{Limitations and Threats to Validity}}
\begin{itemize}[label={},leftmargin=*]
    \item \iconM Discuss potential data leakage effects and their impact on results, and describe how the quality of subjective results was ensured~(\reflimitations).
    % Limitations Internal Validity: "researchers \must discuss potential data leakage effects and their impact on results" (the threat list names evaluation data entering model-improvement pipelines); Limitations Construct Validity: "Researchers \must discuss how they ensured quality of subjective results".
    \item \iconM Transparently report study limitations, including the impact of non-determinism and generalizability constraints~(\reflimitations).
    % Limitations Recommendations / Reliability: "Researchers \must clearly present the limitations of their work without defensiveness or obfuscation"; "Researchers \must discuss the measures taken to increase reliability and reproducibility".
    \item \iconM Specify whether generalization across LLMs or across time was assessed, and discuss model and version differences~(\reflimitations).
    % Limitations External Validity: "Researchers \must discuss the limitations and mitigations of external validity"; covers cross-model transfer and cross-time instability.
    \item \iconM \condition{restricted-sharing} Acknowledge non-disclosed confidential or proprietary components as reproducibility limitations~(\refdesign).
    % Design Recommendations: "Researchers \must clearly describe the tool architecture..., including any dependencies on proprietary tools that affect reproducibility"; Design Challenges: "For commercial tools, researchers \must report all available information and acknowledge unknown aspects as limitations".
    \item \iconS Employ and report strategies to mitigate identified validity and reproducibility threats, such as replication packages, human validation, longitudinal re-runs, triangulation, and sensitivity analysis~(\reflimitations).
    % Limitations Mitigations distributed across threat sections: Replication Packages, Methodological Trustworthiness Measures, Longitudinal Re-Runs, Cost Accounting (Reliability \& Reproducibility); Triangulation, Sensitivity Analysis (External Validity); Human Validation (Construct Validity). No explicit \must; "Researchers \should consider triangulation, reflexivity, audit trails, and peer debriefing as complementary measures".
\end{itemize}


\end{appendices}